% =========================================================================
% Capítulo 3 - Tema 01: Sistemas de Partículas Acopladas
% =========================================================================

\section{Ejercicios de Sistemas de Partículas}

Los siguientes problemas involucran dos o más partículas o cuerpos conectados
mediante varillas rígidas, resortes o vínculos cinemáticos, requiriendo el
análisis del sistema completo: centro de masa, momento lineal y angular
totales, y conservación de energía cuando corresponde.

\subsection{Ejercicios}

% =========================================================================
% Capítulo 3: Cinética de sistemas de partículas
% Tema 01: Sistemas de partículas
% Problema 01: Péndulo invertido sobre bloque móvil
% Origen: Ayudantía 2 (10), 2026
% =========================================================================
\subsubsection{Péndulo invertido sobre bloque móvil}

Un bloque de masa $M$ se desplaza sin fricción a lo largo de una superficie horizontal bajo la acción de una fuerza externa horizontal $F$. Una barra rígida sin masa de longitud $l$ está conectada al bloque mediante una bisagra en el punto $O$, permitiéndole rotar libremente en el plano vertical. En el extremo opuesto de la barra se encuentra una partícula de masa $m$, como se ilustra en la \textbf{Figura \ref{fig:pendulo_bloque}}. Sabiendo que la gravedad actúa hacia abajo y que el sistema parte de una configuración definida por el desplazamiento $x$ del bloque y el ángulo $\theta$ de la barra respecto a la vertical, realice lo siguiente:

\begin{enumerate}
    \item Defina los marcos de referencia y sistemas de coordenadas adecuados.
    \item Analice la cinemática de la partícula y del bloque para obtener sus vectores de posición, velocidad y aceleración.
    \item Analice la cinética del sistema para determinar el sistema de dos ecuaciones diferenciales de movimiento que describen el comportamiento de $x(t)$ y $\theta(t)$.
\end{enumerate}

\begin{figure}[!ht]
    \centering
    \begin{tikzpicture}[>=Latex, scale=1.3]
        % --- Parámetros ---
        \def\ang{35}      % Ángulo theta
        \def\Lbar{2.5}    % Longitud de la barra
        \def\Wblock{1.8}  % Ancho del bloque
        \def\Hblock{0.8}  % Alto del bloque

        % --- Suelo ---
        \fill[pattern=north east lines] (-2.5, -0.2) rectangle (4, 0);
        \draw[thick] (-2.5, 0) -- (4, 0);

        % --- Bloque M ---
        \draw[thick, fill=gray!40] (0, 0.2) rectangle (\Wblock, \Hblock+0.2);
        \node at (\Wblock/2, \Hblock/2+0.2) {$M$};
        
        % Ruedas
        \draw[fill=white] (0.4, 0.1) circle (0.1);
        \draw[fill=white] (\Wblock-0.4, 0.1) circle (0.1);

        % Punto O (Bisagra)
        \coordinate (O) at (\Wblock/2, \Hblock+0.2);
        \filldraw[fill=gray!20] (O) circle (2pt) node[above left, xshift=-2pt] {$O$};

        % --- Barra y Partícula m ---
        \begin{scope}[shift={(O)}, rotate=-\ang]
            \draw[line width=4pt, line cap=round] (0,0) -- (0, \Lbar) node[midway, right] {$l$};
            \draw[line width=3pt, gray, line cap=round] (0,0) -- (0, \Lbar);
            \shadedraw[ball color=gray!30] (0, \Lbar) circle (0.2) node[above right=2pt] {$P$} node[font=\footnotesize] {$m$};
            \coordinate (P) at (0, \Lbar);
        \end{scope}
        \filldraw[] (O) circle (0.6pt);

        % --- Fuerzas y Referencias ---
        % Fuerza F
        \draw[->, ultra thick] (-1, \Hblock/2+0.2) -- (0, \Hblock/2+0.2) node[midway, above left, xshift=-5pt] {$\mathbf{F}$};

        % Eje x de referencia
        \draw[|->, thick] (\Wblock/2, \Hblock+2) -- (\Wblock/2 + 1, \Hblock+2) node[midway, above] {$x$};
        
        % Ángulo theta
        \draw[dashed, gray] (O) -- (\Wblock/2, \Hblock+1.5);
        \draw[->] (\Wblock/2, \Hblock+1.0) arc (90:{90-\ang}:0.8);
        \node at (\Wblock/2+0.3, \Hblock+1.3) {$\theta$};

        % Vector Gravedad
        \draw[->, thick] (-1.5, 2.5) -- (-1.5, 1.5) node[midway, left] {$\mathbf{g}$};

    \end{tikzpicture}
    \caption{Diagrama de una partícula de masa $m$ acoplada a un bloque de masa $M$ con movimiento horizontal y rotación libre.}
    \label{fig:pendulo_bloque}
\end{figure}

% --- SOLUCIÓN ---------------------------------------------------------------
\begin{solucion}{Péndulo invertido sobre bloque móvil}
\textbf{Marcos de referencia y Sistemas de Coordenadas}
\begin{itemize}
    \item $\mathcal{F}$: inercial, fijo al suelo, con origen en la posición inicial de la bisagra $O$ $\longrightarrow\{\vec{E}_x;\vec{E}_y;\vec{E}_z\}$, con $\vec{E}_x$ horizontal hacia la derecha, $\vec{E}_y$ vertical hacia arriba y $\vec{E}_z=\vec{E}_x\times\vec{E}_y$ saliendo de la página.
    \item $\mathcal{A}$: solidario a la varilla $OP$ $\longrightarrow\{\vec{e}_r;\vec{e}_\theta;\vec{e}_z\}$, con $\vec{e}_r$ a lo largo de $OP$, $\vec{e}_z=\vec{E}_z$ y $\vec{e}_\theta=\vec{e}_z\times\vec{e}_r$.
\end{itemize}

\begin{center}
\begin{tikzpicture}[>=Latex, scale=1.3]
    % --- PARÁMETROS ---
    \def\ang{35}   % ángulo theta [°]
    \def\len{1.5}  % largo de los vectores [u]

    % --- BASE FIJA ---
    \draw[->, thick, azul] (0,0) -- (\len,0) node[right] {$\vec{E}_x$};
    \draw[->, thick, azul] (0,0) -- (0,\len) node[above] {$\vec{E}_y$};
    \draw[dashed, gray] (0,0) -- (0,-0.3);

    % --- BASE MÓVIL ---
    \draw[->, thick, myred] (0,0) -- ({90-\ang}:\len) node[above right] {$\vec{e}_r$};
    \draw[->, thick, myred] (0,0) -- ({180-\ang}:\len) node[above left] {$\vec{e}_\theta$};

    % --- ÁNGULO ---
    \draw[->] (0,0.8) arc (90:{90-\ang}:0.8);
    \node at ({90-\ang/2}:1.05) {$\theta$};
    \filldraw (0,0) circle (1.2pt) node[below left] {$O$};
\end{tikzpicture}
\end{center}

De la geometría de la Figura (donde $\theta$ se mide en sentido horario desde la vertical superior $\vec{E}_y$), la relación entre las bases y la velocidad angular de la varilla son:
\begin{equation*}
\begin{split}
    \vec{e}_r &= \sin\theta\,\vec{E}_x + \cos\theta\,\vec{E}_y\\
    \vec{e}_\theta &= -\cos\theta\,\vec{E}_x + \sin\theta\,\vec{E}_y
\end{split}
\end{equation*}
$$\velang{F}{A} = -\dot{\theta}\,\vec{E}_z$$

\subsubsection*{Cinemática del sistema de partículas}
\textbf{Bloque de masa $M$ (punto $O$)}
$$\vec{r}_M = x\,\vec{E}_x \quad\Longrightarrow\quad \marc{\vec{v}_M}{F} = \dot{x}\,\vec{E}_x \quad\Longrightarrow\quad \marc{\vec{a}_M}{F} = \ddot{x}\,\vec{E}_x$$

\textbf{Partícula de masa $m$ (punto $P$)}\\
La posición relativa es $\vec{r}_{P/O} = l\,\vec{e}_r$. Aplicando el Teorema de Transporte dos veces:
$$\marc{\vec{v}_m}{F} = \marc{\vec{v}_M}{F} + \velang{F}{A}\times\vec{r}_{P/O} = \dot{x}\,\vec{E}_x + (-\dot{\theta}\,\vec{E}_z) \times (l\,\vec{e}_r) = \dot{x}\,\vec{E}_x - l\dot{\theta}\,\vec{e}_\theta$$
$$\marc{\vec{a}_m}{F} = \marc{\vec{a}_M}{F} + \der{(-l\dot{\theta}\,\vec{e}_\theta)}{A}{} + \velang{F}{A} \times (-l\dot{\theta}\,\vec{e}_\theta) = \ddot{x}\,\vec{E}_x - l\ddot{\theta}\,\vec{e}_\theta - l\dot{\theta}^2\,\vec{e}_r$$
Sustituyendo $\vec{e}_r$ y $\vec{e}_\theta$ en la base inercial $\{\vec{E}_x, \vec{E}_y\}$:
$$\marc{\vec{a}_m}{F} = \left(\ddot{x} + l\ddot{\theta}\cos\theta - l\dot{\theta}^2\sin\theta\right)\vec{E}_x + \left(-l\ddot{\theta}\sin\theta - l\dot{\theta}^2\cos\theta\right)\vec{E}_y$$

\subsubsection*{Cinética del sistema de partículas}
\textbf{Primera ecuación: Newton para el conjunto $\{M, m\}$ en $\vec{E}_x$}\\
Para el sistema completo $\{M, m\}$, la fuerza axial interna de la varilla se anula por acción y reacción. La única fuerza externa con componente horizontal es $\vec{F} = F\,\vec{E}_x$. Aplicando $\sum \vec{F}_{\text{ext}} \cdot \vec{E}_x = M a_{M,x} + m a_{m,x}$:
$$F = M\ddot{x} + m\left(\ddot{x} + l\ddot{\theta}\cos\theta - l\dot{\theta}^2\sin\theta\right)$$
\respuesta{$\displaystyle (M + m)\ddot{x} + ml\ddot{\theta}\cos\theta - ml\dot{\theta}^2\sin\theta = F$}

\textbf{Segunda ecuación: Newton para $m$ en dirección transversal $\vec{e}_\theta$}\\
Aislando la partícula $m$, las fuerzas que actúan sobre ella son la reacción axial de la varilla $\vec{R} = R\,\vec{e}_r$ y el peso $-mg\,\vec{E}_y$. Proyectando sobre $\vec{e}_\theta$ eliminamos la fuerza de reacción desconocida $R$:
$$\sum \vec{F}_m \cdot \vec{e}_\theta = m\left(\marc{\vec{a}_m}{F}\cdot \vec{e}_\theta\right)$$
Dado que $(-mg\,\vec{E}_y) \cdot \vec{e}_\theta = -mg\sin\theta$ y $\vec{E}_x \cdot \vec{e}_\theta = -\cos\theta$:
$$-mg\sin\theta = m\left(\ddot{x}\,\vec{E}_x - l\ddot{\theta}\,\vec{e}_\theta - l\dot{\theta}^2\,\vec{e}_r\right) \cdot \vec{e}_\theta = m\left(-\ddot{x}\cos\theta - l\ddot{\theta}\right)$$
Dividiendo por $-m$ y ordenando términos, se obtiene la segunda ecuación diferencial:
\respuesta{$\displaystyle l\ddot{\theta} + \ddot{x}\cos\theta - g\sin\theta = 0$}

\patoojo{Proyecta sobre la dirección perpendicular a la fuerza que no conoces. Aquí la varilla solo empuja en $\vec{e}_r$, así que al proyectar la ecuación de $m$ sobre $\vec{e}_\theta$ desaparece la reacción $R$ y no hay que despejarla.}

\patotip{Fíjate en el signo de $g\sin\theta$: es un péndulo \emph{invertido} ($\theta$ medido desde la vertical hacia arriba), por eso la gravedad desestabiliza y el término aparece como $-g\sin\theta$, con el signo contrario al $+g\sin\theta$ del péndulo común.}
\end{solucion}

% =========================================================================
% Capítulo 3: Cinética de sistemas de partículas
% Tema 01: Sistemas de partículas
% Problema 02: Péndulo de mancuerna en forma de T
% Origen: Ayudantía 29, 2026
% =========================================================================
\subsubsection{Péndulo de mancuerna en forma de T}

Una mancuerna de masa $2m$ y largo $2a$, compuesta por dos partículas de masa $m$ en sus extremos, está soldada a una barra rígida sin masa de largo $l$ en su punto medio $P$, formando una estructura en forma de letra $T$. El sistema es libre de rotar en el plano vertical alrededor de una articulación fija en el punto $O$, como se muestra en la \textbf{Figura \ref{fig:mancuerna_T}}. Sabiendo que la gravedad actúa hacia abajo y que el ángulo $\theta$ se mide desde la vertical descendente hacia la barra de longitud $l$, realice lo siguiente:

\begin{enumerate}
    \item Defina los marcos de referencia y sistemas de coordenadas adecuados.
    \item Analice la cinemática del cuerpo rígido, determinando el momento de inercia del sistema respecto al punto de rotación $O$.
    \item Analice la cinética del sistema para obtener la ecuación diferencial de movimiento en términos de $\theta$ utilizando balance de momentum angular y el teorema de trabajo y energía.
\end{enumerate}

\begin{figure}[!ht]
    \centering
    \begin{tikzpicture}[>=Latex, scale=1.3]
        % --- Parámetros ---
        \def\ang{30}      % Ángulo theta
        \def\Lbar{3.0}    % Longitud l
        \def\Aarm{1.2}    % Mitad del largo de la mancuerna (a)

        % --- Soporte Fijo O ---
        \fill[pattern=north east lines] (-0.5, 0.1) rectangle (0.5, 0.4);
        \draw[thick] (-0.5, 0.1) -- (0.5, 0.1);
        \draw[thick, fill=gray!20] (-0.2,0.1) arc (0:180:-0.2) -- cycle;
        
        \coordinate (O) at (0,0);
        \filldraw (O) circle (2pt) node[left, xshift=-5pt] {$O$};

        % --- Barra y Mancuerna ---
        \begin{scope}[rotate=-\ang]
            % Barra principal l
            \draw[line width=4pt] (0,0) -- (0, -\Lbar);
            \draw[line width=4pt] (-\Aarm, -\Lbar) -- (\Aarm, -\Lbar);
            \draw[line width=3pt, gray!40, line cap=round] (0,0) -- (0, -\Lbar);
            \filldraw (O) circle (0.6pt);
            %\node[left] at (0, -\Lbar/2) {$l$};
            
            \coordinate (P) at (0, -\Lbar);
            \node[below left] at (P) {$P$};

            % Travesaño de la mancuerna (2a)
            \draw[line width=3pt, gray!40] (-\Aarm, -\Lbar) -- (\Aarm, -\Lbar);
            \draw[thick] (0,0) -- (0, -\Lbar);
            \draw[thick] (-\Aarm, -\Lbar) -- (\Aarm, -\Lbar);
            \filldraw (P) circle (0.6pt);
            
            % Masas m
            \shadedraw[ball color=gray!30] (-\Aarm, -\Lbar) circle (0.2) node[font=\footnotesize] {$m$};
            \shadedraw[ball color=gray!30] (\Aarm, -\Lbar) circle (0.2) node[font=\footnotesize] {$m$};
            
            % Cota 2a
            \draw[|<->|] (-\Aarm, -\Lbar-0.5) -- (\Aarm, -\Lbar-0.5) node[midway, below] {$2a$};
            
            % Cota l (indicador lateral)
            \draw[|<->|] (-0.8, 0) -- (-0.8, -\Lbar) node[midway, above left] {$l$};
        \end{scope}

        % --- Referencias ---
        % Línea vertical de referencia
        \draw[dashed, gray] (O) -- (0, -2.5);
        
        % Ángulo theta
        \draw[->] (0, -1.0) arc (-90:{-90-\ang}:1.0);
        \node at (-0.3, -1.3) {$\theta$};

        % Vector Gravedad
        \draw[->, thick] (1.5, -0.5) -- (1.5, -1.5) node[midway, right] {$\mathbf{g}$};

    \end{tikzpicture}
    \caption{Esquema de una mancuerna rígida en rotación suspendida de un punto fijo $O$.}
    \label{fig:mancuerna_T}
\end{figure}

% --- SOLUCIÓN ---------------------------------------------------------------
\begin{solucion}{Péndulo de mancuerna en forma de T (por momento angular)}
\textbf{Marcos de referencia y Sistemas de Coordenadas}
\begin{itemize}
    \item $\mathcal{F}$: inercial $\longrightarrow\{\vec{e}_x;\vec{e}_y;\vec{e}_z\}$
    \item $\mathcal{A}$: ligado a la barra, con rotación en $\vec{e}_z$ $\longrightarrow\{\vec{e}_r;\vec{e}_\theta;\vec{e}_z\}$
\end{itemize}

\begin{center}
\begin{tikzpicture}[>=Latex, scale=1.2]
    % --- PARÁMETROS ---
    \def\ang{30}   % ángulo theta [°]
    \def\len{1.6}  % largo de los vectores [u]

    % --- BASE FIJA (e_x hacia abajo, e_y hacia la izquierda, e_z entrando) ---
    \draw[->, thick, azul] (0,0) -- (0,-\len) node[below] {$\vec{e}_x$};
    \draw[->, thick, azul] (0,0) -- (-\len,0) node[left] {$\vec{e}_y$};
    \draw[gray] (0,0) circle (0.12);
    \draw[gray] (-0.085,-0.085) -- (0.085,0.085) (-0.085,0.085) -- (0.085,-0.085);
    \node[gray, above right] at (0.05,0.12) {$\vec{e}_z$};

    % --- BASE MÓVIL ---
    \draw[->, thick, myred] (0,0) -- ({-90-\ang}:\len) node[below left] {$\vec{e}_r$};
    \draw[->, thick, myred] (0,0) -- ({180-\ang}:\len) node[above left] {$\vec{e}_\theta$};

    % --- ÁNGULO ---
    \draw[->] (0,-0.8) arc (-90:{-90-\ang}:0.8);
    \node at ({-90-\ang/2}:1.05) {$\theta$};
    \node[below right] at (0.05,-0.05) {$O$};
\end{tikzpicture}
\end{center}

Relaciones cinemáticas en coordenadas polares (en el plano $x\!-\!y$):
$$\dot{\vec{e}}_r=\dot\theta\,\vec{e}_\theta,\qquad \dot{\vec{e}}_\theta=-\dot\theta\,\vec{e}_r,\qquad \vec{e}_x=\cos\theta\,\vec{e}_r-\sin\theta\,\vec{e}_\theta .$$
Geometría de las masas (dos partículas de masa $m$):
$$\vec{r}_1= \ell\,\vec{e}_r - a\,\vec{e}_\theta, \qquad \vec{r}_2= \ell\,\vec{e}_r + a\,\vec{e}_\theta .$$

\subsubsection*{Análisis cinemático}
Velocidades (derivando en $\mathcal{F}$ y usando $\dot{\vec{e}}_r,\dot{\vec{e}}_\theta$):
$$\begin{aligned}
\vec{v}_1 &= \frac{d}{dt}\big(\ell\,\vec{e}_r - a\,\vec{e}_\theta\big) = \ell\dot\theta\,\vec{e}_\theta + a\dot\theta\,\vec{e}_r,\\[2mm]
\vec{v}_2 &= \frac{d}{dt}\big(\ell\,\vec{e}_r + a\,\vec{e}_\theta\big) = \ell\dot\theta\,\vec{e}_\theta - a\dot\theta\,\vec{e}_r.
\end{aligned}$$

\subsubsection*{Análisis cinético}
Momento angular respecto de $O$ en el marco $\mathcal{F}$:
$$\begin{aligned}
\marc{\vec{H}_O}{F}{}
&= \sum_{i=1}^{2} \vec{r}_i \times m\,\vec{v}_i \\[1mm]
&= (\ell\,\vec{e}_r - a\,\vec{e}_\theta)\times m(\ell\dot\theta\,\vec{e}_\theta + a\dot\theta\,\vec{e}_r)
 \;+\; (\ell\,\vec{e}_r + a\,\vec{e}_\theta)\times m(\ell\dot\theta\,\vec{e}_\theta - a\dot\theta\,\vec{e}_r)\\[1mm]
&= m\dot\theta\Big[\ell^2\,\vec{e}_r\!\times\!\vec{e}_\theta + a^2\,\vec{e}_\theta\!\times\!\vec{e}_r\Big]
 + m\dot\theta\Big[\ell^2\,\vec{e}_r\!\times\!\vec{e}_\theta + a^2\,\vec{e}_\theta\!\times\!\vec{e}_r\Big] \\[1mm]
&= 2m\dot\theta\,(\ell^2+a^2)\,\vec{e}_z .
\end{aligned}$$
De aquí se lee directamente el momento de inercia respecto de $O$, ya que $\marc{\vec{H}_O}{F}{}=I_O\dot\theta\,\vec{e}_z$:
$$I_O = 2m(\ell^2+a^2)$$
Como $O$ es fijo en el inercial,
$$\der{\marc{\vec{H}_O}{F}{}}{F}{} = 2m\ddot\theta\,(\ell^2+a^2)\,\vec{e}_z.$$

Diagrama de cuerpos libres: las únicas fuerzas externas relevantes son los pesos $m\vec{g}$ de cada partícula. Sus momentos respecto de $O$:
$$\marc{\vec{M}_O}{F}{} = \vec{r}_1\times m\vec{g} + \vec{r}_2\times m\vec{g}.$$
Tomando $\vec{g} = -g\,\vec{e}_x$ y usando $\vec{e}_x=\cos\theta\,\vec{e}_r-\sin\theta\,\vec{e}_\theta$:
$$\begin{aligned}
\vec{r}_1\times m\vec{g}
&= (\ell\,\vec{e}_r - a\,\vec{e}_\theta)\times\big[-mg(\cos\theta\,\vec{e}_r-\sin\theta\,\vec{e}_\theta)\big]
= mg(-\ell\sin\theta + a\cos\theta)\,\vec{e}_z,\\[1mm]
\vec{r}_2\times m\vec{g}
&= (\ell\,\vec{e}_r + a\,\vec{e}_\theta)\times\big[-mg(\cos\theta\,\vec{e}_r-\sin\theta\,\vec{e}_\theta)\big]
= mg(-\ell\sin\theta - a\cos\theta)\,\vec{e}_z.
\end{aligned}$$
Por lo tanto,
$$\marc{\vec{M}_O}{F}{}
= \big[mg(-\ell\sin\theta + a\cos\theta) + mg(-\ell\sin\theta - a\cos\theta)\big]\vec{e}_z
= -2mg\ell\sin\theta\,\vec{e}_z .$$

\subsubsection*{Ecuación de movimiento}
Balance de momento angular en $O$:
$$\der{\marc{\vec{H}_O}{F}{}}{F}{} = \marc{\vec{M}_O}{F}{}
\quad\Longrightarrow\quad
2m\ddot\theta(\ell^2+a^2)\,\vec{e}_z = -2mg\ell\sin\theta\,\vec{e}_z.$$
Cancelando factores comunes y proyectando en $\vec{e}_z$:
\respuesta{$\displaystyle \ddot\theta + \frac{g\ell}{\ell^2+a^2}\,\sin\theta = 0$}

\patotruco{Cuando todo el cuerpo gira alrededor de un punto fijo $O$, el balance de momento angular se reduce a $I_O\ddot\theta=M_O$. Aquí $I_O=\sum m_i d_i^2$ con $d_i^2=\ell^2+a^2$ para ambas masas: nos ahorramos todos los productos cruz.}
\end{solucion}

\begin{solucion}{Péndulo de mancuerna en forma de T (por trabajo y energía)}
\textbf{Balance por teorema de trabajo y energía}\\
Con los mismos marcos y las mismas velocidades $\vec{v}_1,\vec{v}_2$ de la solución anterior, la energía cinética total en $\mathcal{F}$ es
$$\marc{T}{F}{}=\sum_{i=1}^{2}\frac{1}{2}\,m\,\vec{v}_i\!\cdot\!\vec{v}_i .$$
Teorema de trabajo y energía en $\mathcal{F}$: la variación de la energía cinética es la potencia de las fuerzas internas y externas,
$$\frac{d}{dt}\big(\marc{T}{F}{}\big)=\sum_{i=1}^{2}\vec{R}_i\!\cdot\!\vec{v}_i .$$

\textbf{Cálculo de la energía cinética}\\
Con $\vec{v}_1=\ell\dot\theta\,\vec{e}_\theta + a\dot\theta\,\vec{e}_r$ y $\vec{v}_2=\ell\dot\theta\,\vec{e}_\theta - a\dot\theta\,\vec{e}_r$,
$$\begin{aligned}
\marc{T}{F}{}
&= \tfrac12 m\,(\ell\dot\theta\,\vec{e}_\theta + a\dot\theta\,\vec{e}_r)\!\cdot\!(\ell\dot\theta\,\vec{e}_\theta + a\dot\theta\,\vec{e}_r)
   + \tfrac12 m\,(\ell\dot\theta\,\vec{e}_\theta - a\dot\theta\,\vec{e}_r)\!\cdot\!(\ell\dot\theta\,\vec{e}_\theta - a\dot\theta\,\vec{e}_r)\\[1mm]
&= \tfrac12 m\,\dot\theta^2\,[\ell^2 + a^2] + \tfrac12 m\,\dot\theta^2\,[\ell^2 + a^2]
= m\dot\theta^2\,[\ell^2+a^2].
\end{aligned}$$
Por lo tanto,
$$\frac{d}{dt}\,\marc{T}{F}{} = 2\,\dot\theta\,\ddot\theta\,m\,[\ell^2+a^2].$$

\textbf{Cálculo de la potencia de las partículas}\\
Sean $\vec{r}_1=\ell\,\vec{e}_r-a\,\vec{e}_\theta$ y $\vec{r}_2=\ell\,\vec{e}_r+a\,\vec{e}_\theta$. La única fuerza externa que trabaja sobre cada partícula es el peso, con $\vec{g}=-g\,\vec{e}_x$ y $\vec{e}_x=\cos\theta\,\vec{e}_r-\sin\theta\,\vec{e}_\theta$. Así,
$$\sum_{i=1}^{2}\vec{R}_i\!\cdot\!\vec{v}_i
= m g(\cos\theta\,\vec{e}_r-\sin\theta\,\vec{e}_\theta)\!\cdot\!(\ell\dot\theta\,\vec{e}_\theta + a\dot\theta\,\vec{e}_r)
 + m g(\cos\theta\,\vec{e}_r-\sin\theta\,\vec{e}_\theta)\!\cdot\!(\ell\dot\theta\,\vec{e}_\theta - a\dot\theta\,\vec{e}_r).$$
Evaluando los productos escalares:
$$\sum_{i=1}^{2}\vec{R}_i\!\cdot\!\vec{v}_i
= mg\,[a\dot\theta\cos\theta-\ell\dot\theta\sin\theta]
 + mg\,[-a\dot\theta\cos\theta-\ell\dot\theta\sin\theta]
= 2mg\,\ell\,\dot\theta\,\sin\theta .$$

\textbf{Igualando}\\
Aplicando el teorema de trabajo y energía:
$$2\,\dot\theta\,\ddot\theta\,m\,[\ell^2+a^2] \;=\; -\,2\,m g\,\ell\,\dot\theta\,\sin\theta.$$
Si $\dot\theta\neq 0$, se simplifica y queda la ecuación de movimiento:
\respuesta{$\displaystyle \ddot\theta+\frac{g\,\ell}{\ell^2+a^2}\,\sin\theta=0$}

\patoojo{Este método solo da \emph{una} ecuación escalar, que es justo lo que necesitamos porque el sistema tiene un único grado de libertad. Pero ojo con el paso final: dividimos por $\dot\theta$, así que la ecuación resultante vale también cuando $\dot\theta=0$ solo porque coincide con la del balance de momento angular.}
\end{solucion}

% =========================================================================
% Capítulo 3: Cinética de sistemas de partículas
% Tema 01: Sistemas de partículas
% Problema 03: Impacto inelástico en péndulo acoplado a carro
% Origen: Ayudantía 29, 2026
% =========================================================================
\subsubsection{Impacto inelástico en péndulo acoplado a carro}

Una partícula $A$ de masa $m/2$ se desplaza con una velocidad inicial $v_0$ hacia la derecha e impacta a una partícula $B$ de masa $m/2$, inicialmente en reposo. La partícula $B$ está conectada a un carro $C$ de masa $m$ mediante una barra rígida sin masa de longitud $l$. El carro se desliza sobre una superficie horizontal sin fricción. Sabiendo que la colisión es perfectamente inelástica, realice lo siguiente:

\begin{enumerate}
    \item Defina los marcos de referencia y sistemas de coordenadas adecuados.
    \item Analice la cinemática del sistema para determinar la velocidad de cada partícula y la velocidad del centro de masa inmediatamente después del impacto.
    \item Determine el máximo ángulo $\theta$ que alcanza la barra después de la colisión utilizando principios de conservación.
\end{enumerate}

\begin{figure}[!ht]
    \centering
    \begin{tikzpicture}[>=Latex, scale=1.4]
        % --- Parámetros ---
        \def\Lbar{2.0}
        \def\Wcar{2.2}
        \def\Hcar{0.6}

        % --- Suelo ---
        \fill[pattern=north east lines] (-3, -0.2) rectangle (3, 0);
        \draw[thick] (-3, 0) -- (3, 0);
        \draw[thin] (-1.0, 0.4) -- (-0.6, 0.1);

        % --- Carro C ---
        \draw[thick, fill=gray!10] (-\Wcar/2, 0.1) rectangle (\Wcar/2, \Hcar+0.1);
        \node at (0.7, \Hcar/2+0.1) {$C$};
        
        \coordinate (O) at (0, \Hcar-0.1);
        \filldraw (O) circle (1.5pt) node[left, xshift=-2pt] {$O$};

        % --- Barra y Partícula B ---
        \draw[line width = 4pt] (O) -- (0, \Hcar-0.1-\Lbar);
        \draw[line width = 3pt, gray!80, line cap = round] (O) -- (0, \Hcar-0.1-\Lbar);
        \shadedraw[ball color=gray!30] (0, \Hcar-0.1-\Lbar) circle (0.18) node[right, xshift=5pt] {$B$};
        \node[left] at (0, \Hcar-0.1-\Lbar/2) {$l$};
        \filldraw (O) circle (0.6pt);
        
        % --- Partícula A ---
        \coordinate (A) at (-2.0, \Hcar-0.1-\Lbar);
        \shadedraw[ball color=gray!30] (A) circle (0.18) node[left, xshift=-5pt] {$A$};
        \draw[->, thick] (A) -- ++(0.7, 0) node[midway, above] {$\mathbf{v}_0$};
        \node[above left, font=\small] at (-0.8, 0.8) {Sin Fricción};

        % --- Gravedad ---
        \draw[->, thick] (2.2, 1.2) -- (2.2, 0.4) node[midway, right] {$\mathbf{g}$};

    \end{tikzpicture}
    \caption{Sistema de impacto inelástico entre una partícula proyectil y un péndulo montado sobre un carro móvil.}
    \label{fig:impacto_carro_pendulo}
\end{figure}


% --- SOLUCIÓN ---------------------------------------------------------------
\begin{solucion}{Impacto inelástico en péndulo acoplado a carro}
\textbf{Planteamiento.} Marco inercial $F$ con base $\{\vec{e}_x,\vec{e}_y,\vec{e}_z\}$ y marco $A$ ligado a la barra $OB$ con base $\{\vec{e}_r,\vec{e}_{\theta},\vec{e}_z\}$. Coordenadas: $x$ (posición horizontal de $O$, medida desde el lugar del choque) y $\theta$ (ángulo de la barra respecto de $\vec{e}_x$). Todas las velocidades se miden en $F$. Se supone que la barra sin masa sólo transmite fuerza en su propia dirección y que, tras el choque, se estudia el sistema $B+C$ (con centro de masa $G$); $A$ interviene sólo mediante el impulso $\hat P$ que le entrega a $B$.
 
\begin{center}
\begin{tikzpicture}[>=Stealth, scale=1.2, font=\small]
  \def\L{2.0}\def\hc{0.9}\def\wc{1.6}\def\yO{0.45}\def\xo{1.4}\def\th{30}
  \pgfmathsetmacro{\xB}{\xo+\L*sin(\th)}
  \pgfmathsetmacro{\yB}{\yO-\L*cos(\th)}
  \pgfmathsetmacro{\yc}{\hc/2}
  \coordinate (O) at (\xo,\yO);
  \coordinate (B) at (\xB,\yB);
  % suelo y carro
  \fill[pattern=north east lines, pattern color=black!55] (-2.2,-0.22) rectangle (4.8,0);
  \draw[thick] (-2.2,0) -- (4.8,0);
  \draw[thick, fill=black!8] (\xo-\wc,0) rectangle (\xo+\wc,\hc);
  \node at (\xo-1.05,\yc) {$m$};
  \node at (\xo+1.15,\yc) {$C$};
  % cota de x
  \draw[dashed, gray] (0,1.5) -- (0,\yO-\L-0.3);
  \draw[dashed, gray] (\xo,\hc) -- (\xo,1.5);
  \draw[<->, gray] (0,1.3) -- (\xo,1.3) node[midway, above] {$x$};
  % barra, vertical de referencia y ángulo
  \draw[dashed, gray] (O) -- ++(0,-\L);
  \draw[->] (B) -- ++(-60:0.9) node[below right] {$\vec{e}_r$};
  \draw[->] (B) -- ++(30:0.9) node[above right] {$\vec{e}_{\theta}$};
  \draw[line width=2.2pt, black!70] (O) -- (B) node[midway, right=2pt] {$l$};
  \fill (O) circle (0.045);
  \node[above left=0pt] at (O) {$O$};
  \draw[<->, gray] ($(O)+(0,-0.9)$) arc[start angle=-90, end angle=-60, radius=0.9];
  \node at ($(O)+(-75:1.25)$) {$\theta$};
  \node[circle, draw, thick, fill=black!8, inner sep=0pt, minimum size=5.5mm] at (B) {$B$};
  % marco F
  \coordinate (F0) at (-1.8,-0.7);
  \draw[->] (F0) -- ++(0.8,0) node[right] {$\vec{e}_y$};
  \draw[->] (F0) -- ++(0,-0.8) node[below] {$\vec{e}_x$};
  \draw (F0) circle (0.07); \fill (F0) circle (0.02);
  \node[above left=-1pt] at (F0) {$\vec{e}_z$};
  \node at ($(F0)+(-0.45,-0.5)$) {$F$};
\end{tikzpicture}
\end{center}
 
\medskip
\textbf{1) Velocidades justo después del choque.}
Choque perfectamente inelástico ($e=0$), con $n$ la dirección normal al plano de colisión ($\vec{e}_y$):
$$e=\frac{v'_{B_n}-v'_{A_n}}{v_{A_n}-v_{B_n}}=0\quad\Rightarrow\quad v'_{B_n}=v'_{A_n}\equiv v'_n .$$
La barra sólo transmite fuerza según $\vec{e}_x$, así que en $\vec{e}_y$ el único impulso sobre $B$ es el de $A$ y se conserva el momentum del par $A$--$B$:
$$m_A\,\vec{v}_{A_n}+\underbrace{m_B\,\vec{v}_{B_n}}_{0}=m_A\,\vec{v}\,'_{A_n}+m_B\,\vec{v}\,'_{B_n}$$
$$\frac m2\,(v_0\vec{e}_y)=\frac m2\,\vec{v}\,'_n+\frac m2\,\vec{v}\,'_n
\;\Rightarrow\; \vec{v}\,'_A=\vec{v}\,'_B=\frac{v_0}{2}\,\vec{e}_y$$
 
\textbf{Impulso sobre $B$.} Como $B$ parte del reposo:
$$\hat P=m_B\,\vec{v}\,'_B=\frac m2\cdot\frac{v_0}{2}\,\vec{e}_y=\frac{m v_0}{4}\,\vec{e}_y .$$
Este impulso produce un movimiento lineal del centro de masa y otro rotacional de la barra.
 
\textbf{Centro de masa de $B+C$.} El centro de masa estaba en reposo antes del choque:
$$\underbrace{(m_B+m_C)\,\vec{v}_G}_{0}+\hat P=(m_B+m_C)\,\vec{v}_G'
\;\Rightarrow\;
\vec{v}_G'=\frac{m v_0/4}{m/2+m}=\frac{m v_0/4}{\tfrac32 m}=\frac{v_0}{6}\,\vec{e}_y .$$
 
\textbf{Velocidad del carro y giro de la barra.} Por cinemática, con $\vec{e}_{\theta}=\vec{e}_y$ en $\theta=0$:
$$\vec{v}_C=\dot x\,\vec{e}_y,\qquad \vec{v}_B=\dot x\,\vec{e}_y+l\dot\theta\,\vec{e}_{\theta} .$$
Con el subíndice $i$ para el instante justo después del choque, hay dos condiciones. La del centro de masa:
$$ \vec{v}_G'=\frac{m_C\,\dot x_i+m_B\,(\dot x_i+l\dot\theta_i)}{m_C+m_B}\,\vec{e}_y=\Big[\dot x_i+\frac13\,l\dot\theta_i\Big]\vec{e}_y=\frac{v_0}{6}\,\vec{e}_y
\;\Rightarrow\; 6\,\dot x_i+2\,l\dot\theta_i=v_0 ,$$
y la velocidad de $B$ obtenida en el choque, $\dot x_i+l\dot\theta_i=v_0/2$, es decir, $2\,\dot x_i+2\,l\dot\theta_i=v_0$. Resolviendo:
$$\dot x_i=0,\qquad \dot\theta_i=\frac{v_0}{2\,l}.$$
\respuesta{$\displaystyle \vec{v}\,'_A=\vec{v}\,'_B=\frac{v_0}{2}\,\vec{e}_y,\qquad \vec{v}\,'_C=\vec{0},\qquad \vec{v}_G'=\frac{v_0}{6}\,\vec{e}_y$}
 
\medskip
\textbf{2) Máximo ángulo de la barra.}
Tras el choque no actúan fuerzas no conservativas (sin roce, la normal no trabaja y la barra rígida sin masa no hace trabajo neto), luego se conserva la energía de $B+C$: ${}^{F}T_f+{}^{F}U_f={}^{F}T_i+{}^{F}U_i$, con estado inicial justo después del choque y final en $\theta_{\max}$.
 
\textbf{Energía potencial} (referencia $U=0$ a la altura de $O$; la del carro es constante):
$${}^{F}U_i=-\frac m2\,g\,l,\qquad {}^{F}U_f=-\frac m2\,g\,l\cos\theta_{\max}.$$
 
\textbf{Energía cinética (König).} ${}^{F}T={}^{F}T^{(1)}+{}^{F}T^{(2)}$, con
$${}^{F}T^{(1)}=\tfrac12\,(m_B+m_C)\,\vec{v}_G\cdot\vec{v}_G,\qquad
{}^{F}T^{(2)}=\tfrac12\sum_k m_k\,(\vec{v}_k-\vec{v}_G)\cdot(\vec{v}_k-\vec{v}_G).$$
\begin{itemize}
  \item \emph{Final:} en $\theta_{\max}$, $\dot\theta=0$ y las partículas no se mueven respecto del centro de masa, así que ${}^{F}T_f^{(2)}=0$.
  \item \emph{Centro de masa:} no hay fuerzas horizontales sobre $B+C$, luego $\vec{v}_{G,f}\cdot\vec{e}_y=\vec{v}_{G,i}\cdot\vec{e}_y=v_0/6$ y ${}^{F}T_i^{(1)}={}^{F}T_f^{(1)}$.
  \item \emph{Inicial:} con $\vec{v}_{C,i}=\vec{0}$, $\vec{v}_{B,i}=\tfrac{v_0}{2}\vec{e}_y$, $\vec{v}_{G,i}=\tfrac{v_0}{6}\vec{e}_y$:
\end{itemize}
$${}^{F}T_i^{(2)}=\frac12\Big[m\Big(0-\frac{v_0}{6}\Big)^2+\frac m2\Big(\frac{v_0}{2}-\frac{v_0}{6}\Big)^2\Big]
=\frac12\Big[\frac{m v_0^2}{36}+\frac{m v_0^2}{18}\Big]=\frac{m v_0^2}{24}.$$
 
\textbf{Balance.} Como ${}^{F}T^{(1)}$ no cambia y ${}^{F}T_f^{(2)}=0$, queda $U_f={}^{F}T_i^{(2)}+U_i$:
$$-\frac m2\,g\,l\cos\theta_{\max}=\frac{m v_0^2}{24}-\frac m2\,g\,l
\;\Rightarrow\;
\cos\theta_{\max}=1-\frac{v_0^2}{12\,g\,l}$$
\respuesta{$\displaystyle \theta_{\max}=\arccos\!\left(1-\frac{v_0^2}{12\,g\,l}\right)$}
Válido mientras el argumento sea mayor que $-1$, es decir, $v_0^2<24\,g\,l$.
 
\medskip
\textbf{Verificación.} (i) Momentum total en $\vec{e}_y$: antes, $m_A v_0=\tfrac12 m v_0$; después, $\tfrac m2\tfrac{v_0}{2}+\tfrac m2\tfrac{v_0}{2}+0=\tfrac12 m v_0$. (ii) Sin König: en $\theta_{\max}$, $B$ y $C$ van a la misma velocidad $u$ y $\tfrac32 m\,u=\tfrac m2\tfrac{v_0}{2}$, así $u=v_0/6$; la energía da $\tfrac{m v_0^2}{16}=\tfrac{m v_0^2}{48}+\tfrac m2 g l(1-\cos\theta_{\max})$, que lleva al mismo resultado. (iii) Si $v_0\to0$, $\theta_{\max}\to0$.
 

\patoojo{Durante el choque, que dura un instante, la barra solo puede transmitir fuerza a lo largo de su propia dirección. Como el impulso de $A$ sobre $B$ es perpendicular a la barra, el carro $C$ no recibe impulso horizontal y por eso sale del choque con $\dot{x}=0$. Toda la velocidad de $B$ se la lleva la rotación de la barra.}

\patotruco{Para el ángulo máximo no hace falta resolver ninguna ecuación diferencial: en $\theta_{\max}$ la barra deja de girar y $B$ y $C$ se mueven juntos con la misma velocidad $u$. Conservación de momentum horizontal ($u=v_0/6$) más conservación de energía, y listo.}
\end{solucion}

% =========================================================================
% Capítulo 3: Cinética de sistemas de partículas
% Tema 01: Sistemas de partículas
% Problema 04: Sistema de correderas ortogonales acopladas
% =========================================================================
\subsubsection{Sistema de correderas ortogonales acopladas}

Una partícula de masa $m_A$ se desliza sin fricción a lo largo de una barra rígida vertical fija. Esta partícula está unida, a través de un brazo rígido y sin masa de longitud $l$, a una segunda partícula de masa $m_B$, la cual se desliza sin fricción a lo largo de una barra rígida horizontal fija. Suponiendo que $\theta$ es el ángulo formado entre la barra horizontal y el brazo, y que la gravedad actúa hacia abajo, realice lo siguiente:

\begin{enumerate}
    \item Defina los marcos de referencia y sistemas de coordenadas adecuados.
    \item Analice la cinemática de las partículas $A$ y $B$, determinando sus vectores de posición y velocidad en función de $\theta$.
    \item Analice la cinética del sistema para determinar la ecuación diferencial del movimiento en términos del ángulo $\theta$.
\end{enumerate}

\begin{figure}[!ht]
    \centering
    \begin{tikzpicture}[>=Latex, scale=1.5]
        % --- Parámetros ---
        \def\ang{40}      % Ángulo theta
        \def\Lrod{3.5}    % Longitud de la barra l
        \def\RailLen{4.0} % Longitud de los rieles

        % Coordenadas
        \coordinate (O) at (0,0);
        % Posición de B (Horizontal)
        \coordinate (B) at ({\Lrod*cos(\ang)}, 0);
        % Posición de A (Vertical) - Nota: La barra baja, coordenada y negativa
        \coordinate (A) at (0, {-\Lrod*sin(\ang)});

        % --- 1. Rieles y Soportes ---
        % Riel Horizontal
        \draw[line width=4pt] (0, -\RailLen)--(0,0) -- (\RailLen, 0);
        \draw[line width=3pt, gray!40] (0, -\RailLen)--(0,0) -- (\RailLen, 0);
        \draw[] (0, -\RailLen)--(0,0) -- (\RailLen, 0);
        % Soporte pared derecha
        \fill[pattern=north east lines] (\RailLen, -0.2) rectangle (\RailLen+0.2, 0.2);
        \draw[thick] (\RailLen, -0.2) -- (\RailLen, 0.2);

        % Riel Vertical
        %\draw[line width=3pt, gray!40] (0,0) -- (0, -\RailLen);
        %\draw[thick] (0,0) -- (0, -\RailLen);
        % Soporte suelo inferior
        \fill[pattern=north east lines] (-0.2, -\RailLen) rectangle (0.2, -\RailLen-0.2);
        \draw[thick] (-0.2, -\RailLen) -- (0.2, -\RailLen);
        
        % Esquina O
        \node[below right, xshift=2pt] at (O) {$O$};

        % --- 2. Brazo de conexión ---
        \draw[line width=3pt] (A) -- (B) node[midway, above left] {$l$};
        \draw[line width=2pt, gray!40] (A) -- (B);

        % --- 3. Masas (Correderas) ---
        % Masa B (Horizontal) - Dibujada como elipse achatada
        \filldraw[fill=gray!20, draw=black, thick] (B) ellipse (0.2 and 0.2);
        \node[above left, yshift=5pt] at (B) {$m_B$};
        \node[below right, yshift=-5pt, font=\footnotesize] at (B) {$B$};

        % Masa A (Vertical) - Dibujada como elipse achatada verticalmente
        \filldraw[fill=gray!20, draw=black, thick] (A) ellipse (0.2 and 0.2);
        \node[above left, xshift=-5pt] at (A) {$m_A$};
        \node[below right, xshift=5pt, font=\footnotesize] at (A) {$A$};

        % --- 4. Ángulo theta ---
        % Arco desde la horizontal (izquierda de B) hacia la barra
       % \draw[dashed, gray] (O) -- (B); 
        \draw[->, thick] ($(B)+(-0.8,0)$) arc (180:{180+\ang}:0.8);
        \node at ($(B)+(-1.1, -0.3)$) {$\theta$};

        % --- 5. Gravedad ---
        \draw[->, thick] (3, -2.0) -- (3, -3.0) node[midway, right] {$\mathbf{g}$};

    \end{tikzpicture}
    \caption{Sistema de dos partículas deslizantes conectadas por una barra rígida de longitud $l$.}
    \label{fig:correderas_ortogonales}
\end{figure}
% =========================================================================
% Capítulo 3: Cinética de sistemas de partículas
% Tema 01: Sistemas de partículas
% Problema 05: Mancuerna deslizante en pista circular
% Origen: Ayudantía 28, 2026
% =========================================================================
\subsubsection{Mancuerna deslizante en pista circular}

Una mancuerna está constituida por dos partículas $A$ y $B$, cada una de masa $m$, conectadas mediante una barra rígida sin masa de longitud $2l$. Los extremos de la mancuerna se deslizan sin fricción por el interior de una pista circular fija de radio $R$ en el plano vertical. Sea $\theta$ el ángulo que forma la vertical descendente con la línea que une el centro de la pista $O$ con el centro de la barra $C$. Considerando que la gravedad actúa hacia abajo, realice lo siguiente:

\begin{enumerate}
    \item Defina los marcos de referencia y sistemas de coordenadas adecuados.
    \item Analice la cinemática del sistema, estableciendo la relación geométrica entre la distancia $a$ (distancia $OC$) y las constantes $R$ y $l$.
    \item Analice la cinética del sistema para determinar la ecuación diferencial del movimiento en términos del ángulo $\theta$.
\end{enumerate}

\begin{figure}[!ht]
    \centering
    \begin{tikzpicture}[>=Latex, scale=1.5]
        % --- Parámetros ---
        \def\R{2.5}       % Radio de la pista
        \def\halfL{1.5}   % Mitad de la longitud de la barra (l)
        % Por Pitágoras, la distancia a = sqrt(R^2 - l^2). 
        % Con R=2.5 y l=1.5 -> a = sqrt(6.25 - 2.25) = 2.0
        \def\distA{2.0}   
        \def\ang{30}      % Ángulo theta

        % Coordenada Origen
        \coordinate (O) at (0,0);

        % --- 1. Pista Circular ---
        % Pared exterior (hachurada)
        \draw[thick] (O) circle (\R+0.2);
        \begin{scope}
            \clip (O) circle (\R+0.4);
            \fill[pattern=north east lines, even odd rule] (O) circle (\R+0.2) circle (\R+0.4);
        \end{scope}
        \draw[thick] (O) circle (\R+0.2); % Redibujar borde

        % --- 2. Elementos Rotados (Barra y Radio) ---
        \begin{scope}[rotate=-\ang] % Rotación horaria theta
            
            % Línea OC (distancia a)
            \draw[dashed] (O) -- (0, -\distA) node[midway, left] {$a$};
            \coordinate (C) at (0, -\distA);

            % Barra AB (Perpendicular a OC)
            \draw[line width=4pt] (-\halfL, -\distA) -- (\halfL, -\distA);
            \draw[line width=3pt, gray!40] (-\halfL, -\distA) -- (\halfL, -\distA);
            \draw[thick] (-\halfL, -\distA) -- (\halfL, -\distA);

            % Partículas
            \coordinate (A) at (-\halfL, -\distA);
            \coordinate (B) at (\halfL, -\distA);

            \shadedraw[ball color=gray!30] (A) circle (0.2) node[right, xshift=5pt, yshift=5pt] {$A$};
            \shadedraw[ball color=gray!30] (B) circle (0.2) node[right, xshift=5pt] {$B$};
            \node[font=\footnotesize] at (A) {$m$}; % Opcional
            \node[font=\footnotesize] at (B) {$m$};

            % Radio R (hacia B)
            \draw[->, thin] (O) -- (\R+0.2,0) node[midway, above right] {$R$};
            
            % Cota 2l
            \draw[|<->|] ($(A)+(0,-0.3)$) -- ($(B)+(0,-0.3)$) node[midway, below] {$2l$};
            \filldraw (C) circle (0.6pt) node[above right, xshift=2pt, font=\footnotesize] {$C$};
        \end{scope}

        % --- 3. Referencias y Ángulos ---
        % Eje Vertical
        \draw[dashed, gray] (O) -- (0, -2.5);
        
        % Ángulo theta
        \draw[->] (0, -1.0) arc (-90:{-90-\ang}:1.0);
        \node at (-0.3, -1.3) {$\theta$};

        % Centro O
        \filldraw (O) circle (1.5pt) node[above left] {$O$};

        % Gravedad
        \draw[->, thick] (3.0, 0) -- (3.0, -1.0) node[midway, right] {$\mathbf{g}$};

    \end{tikzpicture}
    \caption{Sistema de mancuerna deslizante dentro de una pista circular vertical.}
    \label{fig:mancuerna_circular}
\end{figure}

% --- SOLUCIÓN ---------------------------------------------------------------
\begin{solucion}{Mancuerna deslizante en pista circular}
Dado que no hay fricción y el sistema está sujeto únicamente a la fuerza de gravedad (que es conservativa) y a las fuerzas normales (que no realizan trabajo al ser perpendiculares a la trayectoria), la energía mecánica total del sistema se conserva ($\frac{dE}{dt} = 0$). Este es el enfoque más directo para obtener la ecuación de movimiento.

\subsubsection*{Geometría del sistema}
El triángulo formado por el origen $O$ y las partículas $A$ y $B$ es isósceles, ya que los lados $OA$ y $OB$ equivalen al radio $R$. El centro de la varilla es el punto $C$, y la distancia desde $A$ a $C$ es $l$.

\begin{center}
\begin{tikzpicture}[>=Latex, scale=1.4]
    % --- PARÁMETROS ---
    \def\R{2.5}    % radio de la pista [u]
    \def\l{1.5}    % mitad de la varilla [u]
    \def\a{2.0}    % distancia OC [u]

    % --- GEOMETRÍA ---
    \coordinate (O) at (0,0);
    \coordinate (C) at (0,-\a);
    \coordinate (A) at (-\l,-\a);
    \coordinate (B) at (\l,-\a);
    \draw[thick] (A) -- (B);
    \draw[dashed] (O) -- (C) node[midway, left] {$a$};
    \draw (O) -- (A) node[midway, above left] {$R$};
    \draw (O) -- (B) node[midway, above right] {$R$};
    \draw (-0.2,-\a) -- (-0.2,-\a+0.2) -- (0,-\a+0.2);

    % --- PUNTOS ---
    \foreach \P/\pos in {A/left, B/right}
        \shadedraw[ball color=gray!30] (\P) circle (0.13) node[\pos=4pt, yshift=-3pt] {$\P$};
    \filldraw (O) circle (1.2pt) node[above] {$O$};
    \filldraw (C) circle (1.2pt) node[below] {$C$};
    \draw[|<->|] ($(A)+(0,-0.45)$) -- ($(C)+(0,-0.45)$) node[midway, below] {$l$};
\end{tikzpicture}
\end{center}

La distancia perpendicular desde el centro $O$ a la varilla $AB$ (segmento $OC$) se denota por $a$. Aplicando el teorema de Pitágoras en el triángulo rectángulo $OCA$:
$$a = \sqrt{R^2 - l^2}$$
Como la varilla es rígida y sus extremos siempre tocan el círculo, la distancia $OC = a$ permanece constante durante todo el movimiento. El conjunto entero (las dos masas y la varilla) rota de manera rígida alrededor del punto $O$.

\subsubsection*{Energía del sistema}
\textbf{Energía cinética ($T$)}\\
El sistema gira como un cuerpo rígido alrededor de $O$ con una velocidad angular $\dot{\theta}$. Como las partículas $A$ y $B$ se encuentran siempre a una distancia $R$ del origen, la magnitud de la velocidad tangencial de ambas es $v_A = v_B = R\dot{\theta}$. La energía cinética total es la suma de las energías cinéticas de ambas partículas:
$$T = \frac{1}{2}m v_A^2 + \frac{1}{2}m v_B^2 = \frac{1}{2}m(R\dot{\theta})^2 + \frac{1}{2}m(R\dot{\theta})^2 = m R^2 \dot{\theta}^2$$

\textbf{Energía potencial gravitacional ($V$)}\\
Fijamos el nivel de referencia gravitacional ($V = 0$) en el eje horizontal que pasa por $O$, y consideramos el eje vertical hacia abajo como positivo. La posición del centro de masa del sistema compuesto por las dos partículas idénticas se ubica exactamente en el punto $C$. La coordenada vertical del centro de masa es $y_C = a\cos\theta$. La energía potencial de ambas masas en conjunto es equivalente a concentrar toda la masa ($2m$) en su centro de masa $C$:
$$V = -(2m)g y_C = -2mg a\cos\theta = -2mg\sqrt{R^2 - l^2}\cos\theta$$

\subsubsection*{Conservación de la energía y ecuación de movimiento}
La energía mecánica total del sistema es $E = T + V$:
$$E = m R^2 \dot{\theta}^2 - 2mg\sqrt{R^2 - l^2}\cos\theta = \text{constante}$$
Para encontrar la ecuación diferencial de movimiento, derivamos $E$ con respecto al tiempo $t$ y la igualamos a cero:
$$\frac{dE}{dt} = 2m R^2 \dot{\theta} \ddot{\theta} - 2mg\sqrt{R^2 - l^2}(-\sin\theta)\dot{\theta} = 0$$
$$2m\dot{\theta} \left( R^2\ddot{\theta} + g\sqrt{R^2 - l^2}\sin\theta \right) = 0$$
Para un movimiento general ($\dot{\theta} \neq 0$), el término entre paréntesis debe ser nulo. Al dividir todo entre $R^2$, obtenemos la ecuación diferencial que describe el movimiento de la mancuerna:
\respuesta{$\displaystyle \ddot{\theta} + \frac{g\sqrt{R^2 - l^2}}{R^2}\sin\theta = 0$}

\patotruco{Fíjate que las normales de la pista no aparecen en ningún momento: son perpendiculares a la velocidad y no trabajan. Cuando el problema pide solo la ecuación de movimiento y no las reacciones, la energía te ahorra proyectar fuerzas.}

\patotip[squareglasses]{Comprobación rápida: si $l\to0$ la mancuerna colapsa en una partícula doble de masa $2m$ que cae por una pista de radio $R$, y la ecuación se reduce a la del péndulo simple $\ddot\theta+\frac{g}{R}\sin\theta=0$.}
\end{solucion}

% =========================================================================
% Capítulo 3: Cinética de sistemas de partículas
% Tema 01: Sistemas de partículas
% Problema 06: Mancuerna en bloque con pista circular
% Origen: Ayudantía 5 (Sólido rígido), 2026
% =========================================================================
\subsubsection{Mancuerna en bloque con pista circular}

Una mancuerna constituida por dos partículas $A$ y $B$, de masa $m$ cada una, se desliza sin fricción a lo largo de una pista circular de radio $R$ cortada en el interior de un bloque de masa $M$. El bloque, a su vez, se desliza sin fricción a lo largo de una superficie horizontal. Sabiendo que $\theta$ es el ángulo formado por la vertical y la línea que une el centro de la pista $O$ con el centro de la barra $C$, y que $x$ representa el desplazamiento horizontal del bloque, realice lo siguiente:

\begin{enumerate}
    \item Defina los marcos de referencia y sistemas de coordenadas adecuados (fijo y móvil).
    \item Analice la cinemática del sistema para expresar la posición y velocidad de las partículas en función de las coordenadas generalizadas $x$ y $\theta$.
    \item Analice la cinética del sistema para determinar las ecuaciones diferenciales del movimiento utilizando las leyes de Newton (balance de momentum lineal y angular).
\end{enumerate}

\begin{figure}[!ht]
    \centering
    \begin{tikzpicture}[>=Latex, scale=1.4]
        % --- Parámetros ---
        \def\R{2.2}       % Radio de la pista
        \def\halfL{1.4}   % Mitad de la longitud de la barra (l)
        \def\distA{1.7}   % Distancia a (ajustada visualmente < R)
        \def\ang{35}      % Ángulo theta
        \def\BlockSize{3.0} % Tamaño medio del bloque

        % Coordenada Origen del bloque (centro del círculo)
        \coordinate (O) at (0,0);

        % --- 1. Suelo ---
        \draw[thick] (-4, -\BlockSize) -- (4, -\BlockSize);
        \fill[pattern=north east lines] (-4, -\BlockSize-0.2) rectangle (4, -\BlockSize);

        % --- 2. Bloque M ---
        % El bloque es un rectángulo con un agujero circular
        \fill[gray!5, even odd rule] 
            (-\BlockSize, -\BlockSize) rectangle (\BlockSize, \BlockSize)
            (0,0) circle (\R);
        \draw[thick] (-\BlockSize, -\BlockSize) rectangle (\BlockSize, \BlockSize);
        \draw[thick] (0,0) circle (\R);
        
        % Hachurado en el borde interior de la pista (estilo técnico)
        \path[pattern=north east lines, pattern color=gray!80, even odd rule] (0,0) circle (\R+0.2) circle (\R+0.4);
        \filldraw[white, draw= black] (0,0) circle (\R+0.2);
        
        % Etiqueta del Bloque
        \node at (-\BlockSize+0.5, \BlockSize-0.5) {$M$};

        % --- 3. Elementos Geométricos (Internos) ---
        \begin{scope}[rotate=-\ang] 
            % Línea OC (distancia a)
            \draw[dashed] (O) -- (0, -\distA) node[midway, left] {$a$};
            \coordinate (C) at (0, -\distA);
            \filldraw (C) circle (1pt) node[below left, font=\footnotesize] {$C$};

            % Barra AB (Perpendicular a OC)
            \draw[line width=3pt, gray!40] (-\halfL, -\distA) -- (\halfL, -\distA);
            \draw[thick] (-\halfL, -\distA) -- (\halfL, -\distA);

            % Partículas
            \coordinate (A) at (-\halfL, -\distA);
            \coordinate (B) at (\halfL, -\distA);

            \shadedraw[ball color=gray!40] (A) circle (0.2) node[above right, xshift=5pt] {$A$};
            \shadedraw[ball color=gray!40] (B) circle (0.2) node[above right, xshift=5pt] {$B$};
            \node[ font=\footnotesize] at (A) {$m$};
            \node[ font=\footnotesize] at (B) {$m$};

            % Radio R (hacia B, auxiliar para mostrar geometría)
            \draw[->, thin, gray] (O) -- (0,\R+0.2) node[midway,below right, black] {$R$};
            
            % Cota 2l
            \draw[|<->|] ($(A)+(0,-0.3)$) -- ($(B)+(0,-0.3)$) node[midway, below] {$2l$};
        \end{scope}

        % --- 4. Referencias y Cotas Externas ---
        % Centro O
        \filldraw (O) circle (1.5pt) node[above left] {$O$};
        
        % Línea vertical referencia para theta
        \draw[dashed, gray] (O) -- (0, -2.0);
        \draw[->] (0, -0.8) arc (-90:{-90-\ang}:0.8);
        \node at (-0.2, -1.1) {$\theta$};

        % Coordenada x del bloque
        \draw[|->, thick] (0.5, -\BlockSize-0.5) -- (1.5, -\BlockSize-0.5) node[midway, above] {$x$};
        
        % Gravedad
        \draw[->, thick] (3.5, 0.5) -- (3.5, -0.5) node[midway, right] {$\mathbf{g}$};

    \end{tikzpicture}
    \caption{Sistema de mancuerna deslizante dentro de un bloque móvil con cavidad circular.}
    \label{fig:mancuerna_bloque_movil}
\end{figure}

% --- SOLUCIÓN ---------------------------------------------------------------
\begin{solucion}{Mancuerna en bloque con pista circular (por Newton)}
\textbf{Marcos de referencia y Sistemas de Coordenadas}\\
Para este caso podemos darnos cuenta de que tenemos dos cuerpos rígidos, el bloque y la mancuerna.
\begin{itemize}
    \item $\mathcal{N}$: suelo (marco inercial) $\longrightarrow\{\vec{e}_x;\vec{e}_y;\vec{e}_z\}$ cartesiano fijo en $\mathcal{N}$.
    \item $\mathcal{A}$: barra / mancuerna $\longrightarrow\{\vec{e}_r;\vec{e}_\theta;\vec{e}_z\}$ cilíndrico, fijo en la barra y usado para $A$ y $B$.
\end{itemize}

\begin{center}
\begin{tikzpicture}[>=Latex, scale=1.2]
    % --- PARÁMETROS ---
    \def\R{1.6}    % radio de la pista [u]
    \def\ang{35}   % ángulo theta [°]
    \def\len{1.0}  % largo de los vectores [u]

    % --- PISTA Y CENTRO ---
    \draw[thick] (0,0) circle (\R);
    \filldraw (0,0) circle (1.2pt) node[above left] {$O$};
    \draw[dashed, gray] (0,0) -- (0,-\R);

    % --- BASE FIJA (e_x hacia abajo, e_y hacia la derecha) ---
    \draw[->, thick, myred] (0,0) -- (0,-\len) node[below left] {$\vec{e}_x$};
    \draw[->, thick, myred] (0,0) -- (\len,0) node[right] {$\vec{e}_y$};

    % --- BASE MÓVIL ---
    \draw[->, thick, azul] (0,0) -- ({-90+\ang}:\len) node[below right] {$\vec{e}_r$};
    \draw[->, thick, azul] (0,0) -- ({\ang}:\len) node[above right] {$\vec{e}_\theta$};
    \draw[dashed] (0,0) -- ({-90+\ang}:\R);

    % --- ÁNGULO ---
    \draw[->] (0,-0.5) arc (-90:{-90+\ang}:0.5);
    \node at ({-90+\ang/2}:0.72) {$\theta$};
\end{tikzpicture}
\end{center}

\subsubsection*{Análisis cinemático}
\textbf{Bloque}
$$\marc{\vec{r}_O}{N}{} = x\,\vec{e}_y, \qquad
\marc{\vec{v}_O}{N}{} = \dot{x}\,\vec{e}_y, \qquad
\marc{\vec{a}_O}{N}{} = \ddot{x}\,\vec{e}_y$$

\textbf{Partícula $A$}
$$\begin{aligned}
\marc{\vec{r}_A}{N}{} &= \marc{\vec{r}_O}{N}{}+\vec{r}_{A/O} = x\,\vec{e}_y + a\,\vec{e}_r - \ell\,\vec{e}_\theta,\\
\marc{\vec{v}_A}{N}{} &= \dot{x}\,\vec{e}_y + a\dot{\theta}\,\vec{e}_\theta - \ell\dot{\theta}\,\vec{e}_r,\\
\marc{\vec{a}_A}{N}{} &= \ddot{x}\,\vec{e}_y + \left(-a\dot{\theta}^2+\ell\ddot{\theta}\right)\vec{e}_r + \left(a\ddot{\theta}+\ell\dot{\theta}^2\right)\vec{e}_\theta.
\end{aligned}$$

\textbf{Partícula $B$}
$$\begin{aligned}
\marc{\vec{r}_B}{N}{} &= x\,\vec{e}_y + a\,\vec{e}_r + \ell\,\vec{e}_\theta,\\
\marc{\vec{v}_B}{N}{} &= \dot{x}\,\vec{e}_y + a\dot{\theta}\,\vec{e}_\theta + \ell\dot{\theta}\,\vec{e}_r,\\
\marc{\vec{a}_B}{N}{} &= \ddot{x}\,\vec{e}_y + \left(-a\dot{\theta}^2-\ell\ddot{\theta}\right)\vec{e}_r + \left(a\ddot{\theta}-\ell\dot{\theta}^2\right)\vec{e}_\theta.
\end{aligned}$$

\textbf{Otra forma de obtener las mismas aceleraciones: matrices de rotación}\\
En el espacio tridimensional, la orientación de un sistema de referencia respecto a un marco inercial fijo $\mathcal{N}$ se describe mediante una matriz de rotación ortogonal $\mathbf{R} \in SO(3)$. Una de sus propiedades fundamentales es que su derivada temporal multiplicada por su transpuesta define la matriz antisimétrica de la velocidad angular:
$$\dot{\mathbf{R}}\mathbf{R}^T = [\boldsymbol{\omega}]_\times, \qquad
[\boldsymbol{\omega}]_\times = \begin{bmatrix} 0 & -\omega_z & \omega_y \\ \omega_z & 0 & -\omega_x \\ -\omega_y & \omega_x & 0 \end{bmatrix}$$
Multiplicando por la derecha por $\mathbf{R}$ y sabiendo que $\mathbf{R}^T\mathbf{R} = \mathbf{I}$, obtenemos $\dot{\mathbf{R}} = [\boldsymbol{\omega}]_\times \mathbf{R}$. Para la segunda derivada diferenciamos con la regla del producto, y con $\boldsymbol{\alpha} = \dot{\boldsymbol{\omega}}$:
$$\ddot{\mathbf{R}} = \left( [\boldsymbol{\alpha}]_\times + [\boldsymbol{\omega}]_\times^2 \right) \mathbf{R}$$
Para una rotación exclusivamente alrededor del eje inercial $z$ con ángulo variable $\theta(t)$, $\boldsymbol{\omega} = [0, 0, \dot{\theta}]^T$ y $\boldsymbol{\alpha} = [0, 0, \ddot{\theta}]^T$, de modo que
$$[\boldsymbol{\omega}]_\times = \begin{bmatrix} 0 & -\dot{\theta} & 0 \\ \dot{\theta} & 0 & 0 \\ 0 & 0 & 0 \end{bmatrix}, \quad
[\boldsymbol{\alpha}]_\times = \begin{bmatrix} 0 & -\ddot{\theta} & 0 \\ \ddot{\theta} & 0 & 0 \\ 0 & 0 & 0 \end{bmatrix}, \quad
[\boldsymbol{\omega}]_\times^2 = \begin{bmatrix} -\dot{\theta}^2 & 0 & 0 \\ 0 & -\dot{\theta}^2 & 0 \\ 0 & 0 & 0 \end{bmatrix}$$
Con $\mathbf{r}_O = x\,\mathbf{e}_y$ y la posición de cada partícula en el marco local del cuerpo ($\mathbf{r}'_{A/O}=[a,-\ell,0]^T$ y $\mathbf{r}'_{B/O}=[a,\ell,0]^T$), para $P=A,B$:
$$\mathbf{r}_P = \mathbf{r}_O + \mathbf{R}\,\mathbf{r}'_{P/O}, \qquad
\mathbf{v}_P = \mathbf{v}_O + [\boldsymbol{\omega}]_\times \mathbf{R}\,\mathbf{r}'_{P/O}, \qquad
\mathbf{a}_P = \mathbf{a}_O + \left( [\boldsymbol{\alpha}]_\times + [\boldsymbol{\omega}]_\times^2 \right) \mathbf{R}\,\mathbf{r}'_{P/O}$$

\subsubsection*{Análisis cinético}
\begin{center}
\begin{tikzpicture}[>=Latex, scale=0.75]
    % --- PARÁMETROS ---
    \def\R{2.2}        % radio de la pista [u]
    \def\halfL{1.4}    % mitad de la varilla [u]
    \def\distA{1.7}    % distancia a [u]
    \def\ang{35}       % ángulo theta [°]
    \def\BlockSize{3.0}% semilado del bloque [u]

    % --- DCL DEL BLOQUE M ---
    \begin{scope}
        \node[font=\bfseries] at (0, \BlockSize + 0.8) {DCL: Bloque $M$};
        \draw[thick, gray!70] (-\BlockSize, -\BlockSize) rectangle (\BlockSize, \BlockSize);
        \draw[thick, gray!70] (0,0) circle (\R);
        \filldraw (0,0) circle (1.5pt) node[above left] {$O$};
        \draw[->, gray, dashed] (0,0) -- (0, -1.5) node[right] {$\vec{e}_x$};
        \draw[->, gray, dashed] (0,0) -- (1.5, 0) node[above] {$\vec{e}_y$};
        \draw[thick, ->, myred] (0,0) -- (0, -2.2) node[right] {$\vec{F}_g^O$};
        \draw[thick, ->, myred] (-0.5, -\BlockSize) -- (-0.5, -\BlockSize+2) node[left] {$\vec{N}_O$};
        \begin{scope}[rotate=-\ang]
            \coordinate (posA) at (-\halfL, -\distA);
            \coordinate (posB) at (\halfL, -\distA);
        \end{scope}
        \draw[thick, ->, myred] (posA) -- ($(posA)!-1.3cm!(0,0)$) node[left] {$-\vec{N}_A$};
        \draw[dashed, gray] (0,0) -- (posA);
        \draw[thick, ->, myred] (posB) -- ($(posB)!-1.3cm!(0,0)$) node[right] {$-\vec{N}_B$};
        \draw[dashed, gray] (0,0) -- (posB);
    \end{scope}

    % --- DCL DE LA PARTÍCULA A ---
    \begin{scope}[xshift=-6cm, yshift=-8cm]
        \node[font=\bfseries] at (0, 2.5) {DCL: Partícula $A$};
        \coordinate (A) at (0,0);
        \shadedraw[ball color=gray!40] (A) circle (0.3) node[above right, xshift=5pt] {$A$};
        \draw[thick, ->, myred] (A) -- (0, -2) node[right] {$\vec{F}_g^A$};
        \begin{scope}[rotate=-\ang]
            \coordinate (dirOA) at (\halfL, \distA);
            \draw[dashed, gray] (A) -- ($(A)!2cm!(dirOA)$) node[above] {hacia $O$};
            \draw[thick, ->, myred] (A) -- ($(A)!1.5cm!(dirOA)$) node[right] {$\vec{N}_A$};
            \draw[dashed, gray] (1.5, 0) -- (-2, 0);
            \draw[thick, ->, myred] (A) -- (-1.8, 0) node[above] {$\vec{R}_A$};
        \end{scope}
    \end{scope}

    % --- DCL DE LA PARTÍCULA B ---
    \begin{scope}[xshift=6cm, yshift=-8cm]
        \node[font=\bfseries] at (0, 2.5) {DCL: Partícula $B$};
        \coordinate (B) at (0,0);
        \shadedraw[ball color=gray!40] (B) circle (0.3) node[above left, xshift=-5pt] {$B$};
        \draw[thick, ->, myred] (B) -- (0, -2) node[right] {$\vec{F}_g^B$};
        \begin{scope}[rotate=-\ang]
            \coordinate (dirOB) at (-\halfL, \distA);
            \draw[dashed, gray] (B) -- ($(B)!2cm!(dirOB)$) node[below right] {hacia $O$};
            \draw[thick, ->, myred] (B) -- ($(B)!1.5cm!(dirOB)$) node[left] {$\vec{N}_B$};
            \draw[dashed, gray] (1.5, 0) -- (-2, 0);
            \draw[thick, ->, myred] (B) -- (1.8, 0) node[above] {$\vec{R}_B$};
        \end{scope}
    \end{scope}
\end{tikzpicture}
\end{center}

\textbf{Conjunto de fuerzas externas}\\
En este caso, al analizar el bloque, tiene 4 fuerzas: el peso, ejercido por la tierra; la normal ejercida por la superficie; y las dos normales que ejercen las masas de la mancuerna.

\textbf{Bloque}
$$\begin{aligned}
\text{a) }&\vec{F}^{\,O}_g = M g\,\vec{e}_x, \qquad
\text{b) } \vec{N}_O = -N_O\,\vec{e}_x,\\
\text{c) }&-\vec{N}_A = N_A\,\frac{a\vec{e}_x-\ell\vec{e}_y}{R}, \qquad
\text{d) } -\vec{N}_B = N_B\,\frac{a\vec{e}_x+\ell\vec{e}_y}{R}.
\end{aligned}$$
Para conocer la dirección de la fuerza normal de la mancuerna nos podemos recordar que es un círculo, por ende tiene una dirección radial; para la partícula usamos la posición de la partícula relativa al punto $O$ normalizada.

\textbf{Partícula $A$}
$$\begin{aligned}
\text{a) }&\vec{F}^{\,A}_g = m g\,\vec{e}_x = m g\cos\theta\,\vec{e}_r - m g\sin\theta\,\vec{e}_\theta,\\
\text{b) }&\vec{N}_A = -N_A\,\frac{a\vec{e}_x-\ell\vec{e}_y}{R},\qquad
\text{c) } \vec{R}_A= R\,\vec{e}_\theta.
\end{aligned}$$
Mismo análisis para la partícula $B$:
$$\begin{aligned}
\text{a) }&\vec{F}^{\,B}_g = m g\,\vec{e}_x = m g\cos\theta\,\vec{e}_r - m g\sin\theta\,\vec{e}_\theta,\\
\text{b) }&\vec{N}_B = -N_B\,\frac{a\vec{e}_x+\ell\vec{e}_y}{R},\qquad
\text{c) } \vec{R}_B= -R\,\vec{e}_\theta.
\end{aligned}$$

Ahora podemos hacer Newton, pero tenemos muchos Newton por hacer. Podemos hacer varios DCL (diagramas de cuerpo libre); ¿cuál tiene ventaja sobre otro? Aquel en el que conozcamos todas las fuerzas sobre ese eje. En este caso el sistema completo tiene 0 fuerzas externas en la dirección $\vec{e}_y$.

\textbf{Balance de fuerzas del CM en $\vec{e}_y$ para $(A+B+\text{bloque})$}
$$0 = \left( m\,\marc{\vec{a}_A}{N}{} + m\,\marc{\vec{a}_B}{N}{} + M\,\marc{\vec{a}_O}{N}{} \right)\!\cdot\!\vec{e}_y,$$
$$0 = (2m+M)\,\ddot{x} - 2ma\,\dot{\theta}^2\sin\theta + 2ma\,\ddot{\theta}\cos\theta,$$
\respuesta{$\displaystyle \ddot{x}=\frac{2ma}{2m+M}\Big(\dot{\theta}^2\sin\theta-\ddot{\theta}\cos\theta\Big)$}

Ahora necesitamos otra ecuación. ¿Cómo la obtenemos? De otro balance de fuerzas, pero en otra dirección o sobre otro sistema. Lamentablemente, en todo otro sistema hay fuerzas que no conocemos; podemos despejarlas y así resolver, pero hay una manera más fácil que en este caso calza por ser circular: si calculamos el momento angular de la mancuerna, las únicas fuerzas son el peso y las normales, y en el balance de momento angular respecto a $O$ la dirección de $\vec{r}_{A\text{ o }B}$ es la misma que la de su normal, por ende su producto cruz es 0. Esto ya viene con muchos años de análisis, pero es bueno que comiencen a trabajar para tener ese entendimiento de cuáles caminos son más fáciles (de todos modos cualquier forma de obtenerlo es válida, solo que es medio complejo). Ahora, si lo hacemos respecto a $O$ tenemos un problema: es un punto que acelera, por ende nos sale el cochino término de la aceleración del punto.

\textbf{Ecuación de torques respecto de $O$ para el sistema $(A+B)$}
$$\der{\marc{\vec{H}_O}{N}{}}{N}{} = 2mR^2\,\ddot{\theta}\,\vec{e}_z,$$
$$\begin{aligned}
\marc{\vec{M}_O}{N}{}
&= (\vec{r}_A-\vec{r}_O)\times\vec{F}^{\,A}_g + (\vec{r}_B-\vec{r}_O)\times\vec{F}^{\,B}_g + (\vec{r}_O-\vec{r})\times 2m\,\marc{\vec{a}_O}{N}{}\\
&= -2mga\sin\theta\,\vec{e}_z - 2ma\,\ddot{x}\cos\theta\,\vec{e}_z,
\end{aligned}$$
$$2mR^2\,\ddot{\theta} = -2mga\sin\theta - 2ma\,\ddot{x}\cos\theta,$$
\respuesta{$\displaystyle \ddot{\theta}=\frac{a}{R^2}\Big[-g\sin\theta-\ddot{x}\cos\theta\Big]$}

\textbf{Derivada del momento angular respecto a un punto acelerado}\\
El momento angular respecto a un punto $Q$ se define como
$$\marc{\vec{H}_Q}{N}{}=(\vec{r}-\vec{r}_Q)\times m(\marc{\vec{v}}{N}{}-\marc{\vec{v}_Q}{N}{})$$
Si derivamos este vector:
$$\der{\marc{\vec{H}_Q}{N}{}}{N}{}=\der{(\vec{r}-\vec{r}_Q)}{N}{}\times m(\marc{\vec{v}}{N}{}-\marc{\vec{v}_Q}{N}{})+(\vec{r}-\vec{r}_Q)\times m\,\der{(\marc{\vec{v}}{N}{}-\marc{\vec{v}_Q}{N}{})}{N}{}$$
El primer término es el producto cruz de un vector consigo mismo, luego es nulo:
$$\der{\marc{\vec{H}_Q}{N}{}}{N}{}=(\vec{r}-\vec{r}_Q)\times m(\marc{\vec{a}}{N}{}-\marc{\vec{a}_Q}{N}{})=(\vec{r}-\vec{r}_Q)\times \vec{F}-(\vec{r}-\vec{r}_Q)\times m\,\marc{\vec{a}_Q}{N}{}$$
$$\der{\marc{\vec{H}_Q}{N}{}}{N}{}=\vec{M}_Q-(\vec{r}-\vec{r}_Q)\times m\,\marc{\vec{a}_Q}{N}{}$$

\subsubsection*{Ecuaciones diferenciales del movimiento}
Eliminando $\ddot{x}$ mediante la expresión anterior:
$$\left(1+\frac{2ma^2}{(2m+M)R^2}\cos^2\theta\right)\ddot{\theta}
= \frac{a}{R^2}\left(-g\sin\theta+\frac{2ma}{2m+M}\dot{\theta}^2\sin\theta\cos\theta\right),$$
por lo que
\respuesta{$\displaystyle \ddot{\theta}
= \frac{a}{R^2}\,
  \frac{-g\sin\theta+\dfrac{2ma}{2m+M}\dot{\theta}^2\sin\theta\cos\theta}
        {1+\dfrac{2ma^2}{(2m+M)R^2}\cos^2\theta}
\qquad\text{y}\qquad
\ddot{x}=\frac{2ma}{2m+M}\Big(\dot{\theta}^2\sin\theta-\ddot{\theta}\cos\theta\Big)$}

\patotruco{Para elegir el sistema de Newton conviene mirar qué fuerzas \emph{no} actúan en cierta dirección. Con todo el conjunto (bloque + mancuerna), en la dirección horizontal no hay fuerzas externas: la primera ecuación sale gratis y sin incógnitas de reacción.}

\patoojo{Cuando tomas momento angular respecto a un punto que acelera (como $O$, que se mueve con el bloque), no basta con $\dot{\vec{H}}_O=\vec{M}_O$: aparece el término extra $-(\vec{r}-\vec{r}_O)\times m\,\vec{a}_O$. Es el que da el $-2ma\ddot{x}\cos\theta$ en la ecuación de torques.}
\end{solucion}

% =========================================================================
% Capítulo 3: Cinética de sistemas de partículas
% Tema 01: Sistemas de partículas
% Problema 07: Barra giratoria en guía semicircular con resortes
% =========================================================================
\subsubsection{Barra giratoria en guía semicircular con resortes}

Un alambre con forma de semicircunferencia de radio $R$ se encuentra fijo en el plano vertical. Sobre este alambre desliza el extremo $P$ de una barra rígida ideal de masa nula, la cual puede girar libremente alrededor de un eje fijo en el centro geométrico $O$. La barra sostiene dos partículas: una de masa $2m$ en el extremo $P$ y otra de masa $m$ en el extremo opuesto $Q$. La distancia entre $O$ y $P$ es $R$, mientras que la distancia entre $O$ y $Q$ es $d$. Además, la partícula $P$ está conectada a dos resortes idénticos de constante elástica $k$ y longitud natural $l_0$, los cuales siguen la curvatura del alambre y están anclados en los extremos de su diámetro horizontal. Considerando que la gravedad actúa hacia abajo y definiendo $\theta$ como el ángulo que forma el segmento $OP$ con la vertical ascendente, realice lo siguiente:

\begin{enumerate}
    \item Defina los marcos de referencia y sistemas de coordenadas adecuados.
    \item Analice la cinemática del sistema para expresar la velocidad de las partículas $P$ y $Q$ en función del ángulo $\theta$.
    \item Analice la cinética del sistema para determinar la ecuación diferencial del movimiento utilizando métodos de energía o fuerzas.
\end{enumerate}

\begin{figure}[!ht]
    \centering
    \begin{tikzpicture}[>=Latex, scale=1.5]
        % --- Parámetros ---
        \def\R{2.2}       % Radio del semicírculo (R)
        \def\distd{1.4}   % Distancia d (OQ)
        \def\ang{25}      % Ángulo theta (desviación de la vertical)
        
        % Coordenadas Clave
        \coordinate (O) at (0,0);
        % P en coordenadas polares (90 - theta)
        \coordinate (P) at ({90-\ang}:\R);
        % Q en coordenadas polares (90 - theta + 180) -> (270 - theta)
        \coordinate (Q) at ({270-\ang}:\distd);
        % Extremos del diámetro
        \coordinate (BaseL) at (-\R, 0);
        \coordinate (BaseR) at (\R, 0);

        % --- 1. Soportes y Línea Base ---
        \draw[dashed, gray] (BaseL) -- (BaseR);
        
        % Soportes en los extremos
        \foreach \x in {-\R, \R} {
            \fill[pattern=north east lines] (\x-0.2, -0.1) rectangle (\x+0.2, 0);
            \draw[thick] (\x-0.2, 0) -- (\x+0.2, 0);
        }

        % --- 2. Resortes en la Guía (Arcos) ---
        % Resorte Izquierdo (desde 180 hasta P)
        \draw[decorate, decoration={coil, aspect=0.4, segment length=3pt, amplitude=2.5pt}] 
            (BaseL) arc (180:{90-\ang}:\R);
        % Resorte Derecho (desde 0 hasta P)
        \draw[decorate, decoration={coil, aspect=0.4, segment length=3pt, amplitude=2.5pt}] 
            (BaseR) arc (0:{90-\ang}:\R);
            
        % Línea guía del alambre (borde fino para definir la forma)
        \draw[thin, gray] (\R,0) arc (0:180:\R);

        % --- 3. Barra y Masas ---
        \draw[line width=2pt, gray!50] (Q) -- (P);
        \draw[thick] (Q) -- (P);

        % Masa P (2m)
        \shadedraw[ball color=gray!30] (P) circle (0.22) node[above right, xshift=5pt] {$P$};
        \node[font=\footnotesize] at (P) {$2m$};

        % Masa Q (m)
        \shadedraw[ball color=gray!30] (Q) circle (0.18) node[below left, xshift =-5pt] {$Q$};
        \node[font=\footnotesize] at (Q) {$m$};

        % --- 4. Cotas y Referencias ---
        % Centro O
        \filldraw (O) circle (1.5pt) node[above right, font=\small, xshift=3pt] {$O$};
        
        % Eje Vertical de Referencia
        \draw[dashed] (0, 0.5) -- (0, \R+0.3);
        
        % Ángulo theta
        \draw[->] (0, 0.8) arc (90:{90-\ang}:0.8);
        \node at (0.25, 1.0) {$\theta$};
        
        % Etiqueta R (Radio)
        \draw[thin, ->] (O) -- (135:\R) node[midway, fill=white, inner sep=1pt] {$R$};
        
        % Etiqueta d
        %\draw[thin, |-|] ($(O)!0.2!(90:\distd)$) -- ($(Q)!0.2!(90:\distd)$); 
        % Ajuste manual de cota d para que sea paralela a la barra
        \draw[thin, |<->|] ($(O)!0.15!({270-\ang+90}:1)$) -- ($(Q)!0.15!({270-\ang+90}:1)$) 
            node[midway, below right] {$d$};

        % Gravedad
        \draw[->, thick] (1.3, 1.0) -- (1.3, 0.1) node[midway, right] {$\vec{g}$};

    \end{tikzpicture}
    \caption{Sistema de barra rotatoria con masas en los extremos y resortes sobre guía semicircular.}
    \label{fig:barra_semicirculo_resortes}
\end{figure}
% =========================================================================
% Capítulo 3: Cinética de sistemas de partículas
% Tema 01: Sistemas de partículas
% Problema 08: Tres péndulos acoplados simétricamente
% =========================================================================
\subsubsection{Tres péndulos acoplados simétricamente}

Tres péndulos simples, todos de longitud $l$, están dispuestos simétricamente suspendidos de un techo horizontal fijo. La masa central, de valor $2m$, está conectada a las masas exteriores, cada una de valor $m$, mediante resortes ideales de constante elástica $k$. El sistema vibra exclusivamente en el plano vertical de la figura. Considerando que la gravedad actúa hacia abajo y asumiendo pequeñas oscilaciones, realice lo siguiente:

\begin{enumerate}
    \item Defina los marcos de referencia y las coordenadas generalizadas para describir el movimiento de las tres masas.
    \item Determine las matrices de inercia (masa) y de rigidez del sistema para obtener las ecuaciones de movimiento en forma matricial.
    \item Calcule las frecuencias naturales (modos normales) y los vectores propios correspondientes que describen los desplazamientos relativos de los péndulos.
\end{enumerate}

\begin{figure}[!ht]
    \centering
    \begin{tikzpicture}[>=Latex, scale=1.3]
        % --- Parámetros ---
        \def\PendL{2.5}   % Longitud l
        \def\Sep{2.2}     % Separación horizontal entre péndulos
        
        % Coordenadas de los pivotes
        \coordinate (P1) at (-\Sep, 0);
        \coordinate (P2) at (0, 0);
        \coordinate (P3) at (\Sep, 0);
        
        % Coordenadas de las masas (Posición de equilibrio)
        \coordinate (M1) at (-\Sep, -\PendL);
        \coordinate (M2) at (0, -\PendL);
        \coordinate (M3) at (\Sep, -\PendL);

        % --- 1. Techo ---
        \fill[pattern=north east lines] (-\Sep-1, 0) rectangle (\Sep+1, 0.2);
        \draw[thick] (-\Sep-1, 0) -- (\Sep+1, 0);

        % --- 2. Varillas de los péndulos ---
        \draw[thick] (P1) -- (M1) node[midway, left] {$l$};
        \draw[thick] (P2) -- (M2) node[midway, left] {$l$};
        \draw[thick] (P3) -- (M3) node[midway, right] {$l$};

        % --- 3. Resortes ---
        % Resorte izquierdo (entre M1 y M2)
        \draw[decoration={coil, aspect=0.4, segment length=4pt, amplitude=3pt, pre length=3pt, post length=3pt}, decorate] 
            (M1) -- (M2);
        \node[above, yshift=3pt] at ($(M1)!0.5!(M2)$) {$k$};

        % Resorte derecho (entre M2 y M3)
        \draw[decoration={coil, aspect=0.4, segment length=4pt, amplitude=3pt, pre length=3pt, post length=3pt}, decorate] 
            (M2) -- (M3);
        \node[above, yshift=3pt] at ($(M2)!0.5!(M3)$) {$k$};

        % --- 4. Masas ---
        % Masa 1 (m)
        \filldraw (P1) circle (1.5pt);
        \shadedraw[ball color=gray!30] (M1) circle (0.2) node[] {$m$};
        
        % Masa 2 (2m) - Un poco más grande
        \filldraw (P2) circle (1.5pt);
        \shadedraw[ball color=gray!30] (M2) circle (0.25) node[] {$2m$};
        
        % Masa 3 (m)
        \filldraw (P3) circle (1.5pt);
        \shadedraw[ball color=gray!30] (M3) circle (0.2) node[] {$m$};

        % --- 5. Gravedad ---
        \draw[->, thick] (\Sep+0.8, -0.5) -- (\Sep+0.8, -1.5) node[midway, right] {$\mathbf{g}$};

    \end{tikzpicture}
    \caption{Sistema de tres péndulos acoplados mediante resortes lineales.}
    \label{fig:tres_pendulos_acoplados}
\end{figure}
% =========================================================================
% Capítulo 3: Cinética de sistemas de partículas
% Tema 01: Sistemas de partículas
% Problema 09: Mecanismo rotatorio con resorte y masa deslizante
% =========================================================================
\subsubsection{Mecanismo rotatorio con resorte y masa deslizante}

Un mecanismo articulado rota alrededor de un eje vertical con una velocidad angular constante $\Omega$. El sistema está compuesto por un disco pequeño $B$ y un collarín deslizante $C$, ambos de masa $m$. Los eslabones rígidos $AB$ y $BC$ tienen longitud $L$ y masa despreciable. El collarín $C$ desliza sobre el eje vertical y está conectado al punto fijo $A$ mediante un resorte lineal de constante de rigidez $k$ y longitud natural $2l$. El resorte se encuentra enrollado alrededor del eje. Sabiendo que $\theta$ es el ángulo que forman los eslabones con la vertical y que la gravedad actúa hacia abajo, realice lo siguiente:

\begin{enumerate}
    \item Defina los marcos de referencia y sistemas de coordenadas adecuados (fijo y rotatorio).
    \item Analice la cinemática del sistema para determinar las expresiones de energía cinética y potencial en función del ángulo $\theta$.
    \item Determine los dos posibles valores de equilibrio constante para $\theta$ en el rango $0 \le \theta < \pi/2$ y analice la estabilidad de dichas posiciones en función de la rigidez $k$ y la velocidad $\Omega$.
\end{enumerate}

\begin{figure}[!ht]
    \centering
    \begin{tikzpicture}[>=Latex, scale=1.4]
        % --- Parámetros ---
        \def\Llink{2.5}   % Longitud L de los eslabones
        \def\ang{35}      % Ángulo theta
        \def\Wslider{0.6} % Ancho del collarín
        \def\Hslider{0.4} % Alto del collarín
        
        % Coordenadas
        \coordinate (A) at (0,0); % Punto fijo inferior
        % Geometría: Triángulo isósceles. Altura total h = 2*L*cos(theta)
        \coordinate (C) at (0, {2*\Llink*cos(\ang)}); % Posición del collarín
        \coordinate (B) at ({\Llink*sin(\ang)}, {\Llink*cos(\ang)}); % Posición de la masa B
        \coordinate (TopSupport) at (0, {2*\Llink*cos(\ang) + 1.0});

        % --- 1. Eje y Soportes ---
        % Indicador de Rotación Omega
        \draw[->, thick] (0.3, {2*\Llink*cos(\ang) + 0.6}) arc (0:255:0.3 and 0.1);
        \node[right] at (0.3, {2*\Llink*cos(\ang) + 0.4}) {$\Omega$};
        % Eje Vertical
        \draw[line width=6pt] (0, -0.5) -- (0, {2*\Llink*cos(\ang) + 1.2});
        \draw[line width=5pt, gray!20] (0, -0.5) -- (0,{2*\Llink*cos(\ang) + 1.2});
        \draw[thick] (0, -0.5) -- (0, {2*\Llink*cos(\ang) + 1.2});

        % Soporte Inferior
        \fill[pattern=north east lines] (-0.5, -0.5) rectangle (0.5, -0.3);
        \draw[thick] (-0.5, -0.3) -- (0.5, -0.3);

        % Soporte Superior
        \fill[pattern=north east lines] (-0.5, {2*\Llink*cos(\ang) + 1.0}) rectangle (0.5, {2*\Llink*cos(\ang) + 1.2});
        \draw[thick] (-0.5, {2*\Llink*cos(\ang) + 1.0}) -- (0.5, {2*\Llink*cos(\ang) + 1.0});


        % --- 2. Resorte (alrededor del eje entre A y C) ---
        \draw[decoration={coil, aspect=0.5, segment length=5pt, amplitude=4pt, pre length=5pt, post length=5pt}, decorate, gray!80, thick] 
            (A) -- (C);
        \node[left, xshift=-5pt] at (0, {\Llink*cos(\ang)}) {$k$};

        % --- 3. Eslabones y Masas ---
        % Eslabón AB
        \draw[line width=4pt] (A) -- (B);
        \draw[line width=3pt, gray!40] (A) -- (B);
        \draw[thick] (A) -- (B) node[midway, below right] {$L$};
        
        % Eslabón CB
        \draw[line width=4pt] (C) -- (B);
        \draw[line width=3pt, gray!40] (C) -- (B);
        \draw[thick] (C) -- (B) node[midway, above right] {$L$};

        % Collarín C
        \shadedraw[color=gray!20, draw=black] ($(C)+(-\Wslider/2, -\Hslider/2)$) rectangle ($(C)+(\Wslider/2, \Hslider/2)$);
        \node[above left, xshift=-5pt, yshift=7pt] at (C) {$C$};
        %\filldraw (C) circle (1.5pt); % Pivote en C
        \node at (C) {$m$};

        % Punto fijo A
        \filldraw (A) circle (1.5pt) node[left, xshift=-5pt] {$A$};

        % Masa B (Disco)
        \shadedraw[ball color=gray!30] (B) circle (0.25) node[right, xshift=5pt] {$B$};
        %\filldraw (B) circle (1.5pt);
        \node[] at (B) {$m$};

        % --- 4. Ángulos theta ---
        % Ángulo en A
        \draw[thin] (0, 0.6) arc (90:{90-\ang}:0.6);
        \node at (0.2, 0.8) {$\theta$};
        
        % Ángulo en C
        \draw[thin] (0, {2*\Llink*cos(\ang)-0.6}) arc (270:{270+\ang}:0.6);
        \node at (0.2, {2*\Llink*cos(\ang)-0.8}) {$\theta$};

    \end{tikzpicture}
    \caption{Mecanismo rotatorio con control elástico y masa centrífuga.}
    \label{fig:gobernador_centrifugo}
\end{figure}
% =========================================================================
% Capítulo 3: Cinética de sistemas de partículas
% Tema 01: Sistemas de partículas
% Problema 10: Péndulo con masa deslizante y resorte
% =========================================================================
\subsubsection{Péndulo con masa deslizante y resorte}

Una partícula de masa $m_1$ está unida al extremo de una barra rígida ideal sin masa de longitud $L$, la cual está articulada en el punto $O$ y puede rotar libremente en el plano vertical. Una segunda masa $m_2$ (cuenta o collarín) puede deslizarse libremente a lo largo de la barra lisa. Esta masa $m_2$ está conectada al pivote $O$ mediante un resorte lineal de constante elástica $k$ y longitud natural $L_0$. Sabiendo que la gravedad actúa hacia abajo, realice lo siguiente:

\begin{enumerate}
    \item Defina un conjunto completo e independiente de coordenadas generalizadas para describir la configuración del sistema.
    \item Analice la cinemática del sistema determinando los vectores de posición y velocidad de ambas masas.
    \item Derive las ecuaciones diferenciales del movimiento aplicando principios de momentum (cantidad de movimiento y momento angular).
\end{enumerate}

\begin{figure}[!ht]
    \centering
    \begin{tikzpicture}[>=Latex, scale=1.4]
        % --- Parámetros ---
        \def\ang{-50}      % Ángulo de la barra
        \def\Lrod{3.5}     % Longitud L
        \def\Rbead{1.8}    % Posición de m2
        
        % Coordenada Origen
        \coordinate (O) at (0,0);

        % --- 1. Techo y Pivote ---
        \fill[pattern=north east lines] (-1.0, 0) rectangle (1.5, 0.3);
        \draw[thick] (-1.0, 0) -- (1.5, 0);
        
        % Pivote O
        \filldraw[thick] (O) circle (3pt);
        \node[below left, yshift=3pt] at (O) {$O$};

        % --- 2. Barra y Resorte ---
        % Barra L
        \draw[line width=5pt] (O) -- (\ang:\Lrod);
        \draw[line width=4pt, gray!30, line cap=round] (O) -- (\ang:\Lrod);
        \draw[thick] (O) -- (\ang:\Lrod);
        \filldraw[thick] (O) circle (0.5pt);

        % Resorte (Desde O hasta m2)
        % Usamos una decoración coil sobre la barra
        \draw[decoration={coil, aspect=0.4, segment length=5pt, amplitude=5pt, pre length=1pt, post length=0pt}, decorate, black] 
            (O) -- (\ang:\Rbead);
        \node[above right, xshift=4pt] at (\ang:{\Rbead/2}) {$k$};

        % --- 3. Masas ---
        % Masa m2 (Cuenta deslizante)
        \coordinate (M2) at (\ang:\Rbead);
        \shadedraw[ball color=gray!30] (M2) circle (0.35) node[] {$m_2$};

        % Masa m1 (Extremo)
        \coordinate (M1) at (\ang:\Lrod);
        \shadedraw[ball color=gray!30] (M1) circle (0.25) node[] {$m_1$};

        % --- 4. Cotas y Gravedad ---
        % Cota L (Longitud total)
        \draw[|<->|] ($(O) + (\ang-90:0.5)$) -- ($(M1) + (\ang-90:0.5)$) node[midway, below left] {$L$};

        % Vector Gravedad
        \draw[->, thick] (2.0, -1.0) -- (2.0, -2.0) node[midway, right] {$g$};

    \end{tikzpicture}
    \caption{Sistema de péndulo compuesto con una masa deslizante sujeta a un resorte.}
    \label{fig:pendulo_masa_deslizante}
\end{figure}
% =========================================================================
% Capítulo 3: Cinética de sistemas de partículas
% Tema 01: Sistemas de partículas
% Problema 11: Sistema de osciladores acoplados con péndulo
% =========================================================================
\subsubsection{Sistema de osciladores acoplados con péndulo}

Considere el sistema mecánico mostrado en la Figura \ref{fig:osciladores_pendulo}. El sistema consta de dos bloques de masas $m_1$ y $m_2$ que se deslizan sobre una superficie horizontal sin fricción. El bloque $m_1$ está conectado a una pared fija mediante un resorte de constante $k_1$ y al bloque $m_2$ mediante un resorte de constante $k_2$. Adicionalmente, una partícula de masa $m_3$ cuelga del bloque $m_2$ mediante una barra rígida sin masa de longitud $L$, formando un péndulo. Asumiendo que la gravedad actúa hacia abajo, realice lo siguiente:

\begin{enumerate}
    \item Defina los marcos de referencia y sistemas de coordenadas adecuados.
    \item Analice la cinemática del sistema describiendo la posición y velocidad de cada masa en función de las coordenadas generalizadas $x_1$, $x_2$ y $\theta$.
    \item Obtenga las ecuaciones diferenciales del movimiento para las coordenadas generalizadas $x_1$, $x_2$ y $\theta$ aplicando principios de momentum (cantidad de movimiento lineal y angular).
\end{enumerate}

\begin{figure}[!ht]
    \centering
    \begin{tikzpicture}[>=Latex, scale=1.4]
        % --- Parámetros ---
        \def\BlockW{1.4}  % Ancho bloque
        \def\BlockH{1.0}  % Alto bloque
        \def\WheelR{0.15} % Radio rueda
        \def\SpringL{1.5} % Largo resorte visual
        \def\PendL{2.0}   % Largo péndulo
        \def\ang{30}      % Ángulo theta
        
        % Coordenadas base (Centros de los bloques)
        \coordinate (Wall) at (0, 0.5*\BlockH + \WheelR);
        \coordinate (M1) at (\SpringL + 0.5*\BlockW, 0.5*\BlockH + \WheelR);
        \coordinate (M2) at (2*\SpringL + 1.5*\BlockW, 0.5*\BlockH + \WheelR);
        \coordinate (Pivot) at ($(M2)+(0,-0.3)$);
        \coordinate (M3) at ($(Pivot) + ({sin(\ang)*\PendL}, {-cos(\ang)*\PendL})$);

        % --- 1. Suelo y Pared ---
        % Pared
        \draw[thick] (0, 0) -- (0, 2.0);
        \fill[pattern=north east lines] (-0.3, 0) rectangle (0, 2.0);
        
        % Suelo
        \draw[thick] (0, 0) -- (7.5, 0);
        \fill[pattern=north east lines] (-0.3, -0.1) rectangle (7.5, 0);

        % --- 2. Resortes ---
        % Resorte k1
        \draw[decoration={coil, aspect=0.4, segment length=4pt, amplitude=4pt, pre length=3pt, post length=3pt}, decorate] 
            (0, 0.5*\BlockH + \WheelR) -- ($(M1) - (0.5*\BlockW, 0)$);
        \node[above, yshift=5pt] at ($(Wall)!0.4!(M1)$) {$k_1$};

        % Resorte k2
        \draw[decoration={coil, aspect=0.4, segment length=4pt, amplitude=4pt, pre length=3pt, post length=3pt}, decorate] 
            ($(M1) + (0.5*\BlockW, 0)$) -- ($(M2) - (0.5*\BlockW, 0)$);
        \node[above, yshift=5pt] at ($(M1)!0.5!(M2)$) {$k_2$};

        % --- 3. Bloques m1 y m2 ---
        % Bloque 1
        \draw[thick, fill=white] ($(M1) - (0.5*\BlockW, 0.5*\BlockH)$) rectangle ($(M1) + (0.5*\BlockW, 0.5*\BlockH)$);
        \node at (M1) {$m_1$};
        % Ruedas Bloque 1
      
        % Corrección posición ruedas:
        \draw[fill=white] ($(M1) + (-0.3*\BlockW, -0.5*\BlockH)$) circle (\WheelR);
        \draw[fill=white] ($(M1) + (0.3*\BlockW, -0.5*\BlockH)$) circle (\WheelR);

        % Bloque 2
        \draw[thick, fill=white] ($(M2) - (0.5*\BlockW, 0.5*\BlockH)$) rectangle ($(M2) + (0.5*\BlockW, 0.5*\BlockH)$);
        \node at (M2) {$m_2$};
        % Ruedas Bloque 2
        \draw[fill=white] ($(M2) + (-0.3*\BlockW, -0.5*\BlockH)$) circle (\WheelR);
        \draw[fill=white] ($(M2) + (0.3*\BlockW, -0.5*\BlockH)$) circle (\WheelR);

        % --- 4. Péndulo m3 ---
        \draw[line width=1.5pt] (Pivot) -- (M3); % Barra sin masa
        \draw[double] (Pivot) -- (M3); 
        \filldraw[fill=white] (Pivot) circle (2pt); % Pivote
        \shadedraw[ball color=white] (M3) circle (0.25) node[] {$m_3$};

        % --- 5. Coordenadas y Vectores ---
        % Coordenada x1
        \draw[->, thick] ($(Wall) + (0, 0.5*\BlockH)$) -- ++(0, 0.6) -- ($(M1)+(0, 0.5*\BlockH+0.6)$) node[midway, above] {$x_1$};
        
        % Coordenada x2
        \draw[->, thick] ($(Wall) + (0, 0.5*\BlockH)$) -- ++(0, 0.2) -- ($(M2)+(0, 0.5*\BlockH+0.2)$) node[midway, above] {$x_2$};

        % Ángulo theta
        \draw[dashed] (Pivot) -- ++(0, -1.5);
        \draw[->] ($(Pivot) + (0, -0.8)$) arc (-90:{-90+\ang}:0.8);
        \node at ($(Pivot) + (0.3, -1.0)$) {$\theta$};

        % Gravedad
        \draw[->, thick] (6.5, 1.5) -- (6.5, 0.5) node[midway, right] {$g$};

    \end{tikzpicture}
    \caption{Sistema de masas acopladas por resortes con un péndulo suspendido.}
    \label{fig:osciladores_pendulo}
\end{figure}
% =========================================================================
% Capítulo 3: Cinética de sistemas de partículas
% Tema 01: Sistemas de partículas
% Problema 12: Péndulo acoplado a soporte con resorte horizontal
% =========================================================================
\subsubsection{Péndulo acoplado a soporte con resorte horizontal}

Se considera un sistema compuesto por un carro de masa $M$ que se desplaza sin fricción sobre una guía horizontal. El carro está unido a una pared fija mediante un resorte lineal de constante elástica $k$ y longitud natural $l_0$. Desde el centro del carro, en el punto $O$, pende una varilla rígida de longitud $L$ y masa despreciable, en cuyo extremo se encuentra una partícula de masa $m$. El movimiento del carro se describe mediante la coordenada horizontal $x(t)$, mientras que la orientación del péndulo se define por el ángulo $\theta(t)$ medido respecto a la vertical, como se muestra en la \textbf{Figura \ref{fig:pendulo_resorte_movil}}. Se asume que la gravedad $g$ actúa en dirección vertical hacia abajo.

\begin{enumerate}
    \item Defina los marcos de referencia y sistemas de coordenadas adecuados para describir la configuración del sistema.
    \item Determine y analice la cinemática de la partícula $m$ y de la masa $M$ en términos de las coordenadas generalizadas $\{x, \theta\}$.
    \item Determine y analice la cinética del sistema para obtener las ecuaciones diferenciales de movimiento.
\end{enumerate}

\begin{figure}[!ht]
    \centering
    \begin{tikzpicture}[>=Latex, scale=1.3]
        % --- CONFIGURACIÓN DE PARÁMETROS ---
        \def\xCar{1.5}      % Posición x del carro
        \def\ang{30}       % Ángulo theta
        \def\L{2.5}        % Longitud del péndulo
        \def\wCar{1.0}     % Ancho del carro
        \def\hCar{0.6}     % Alto del carro

        % --- 1. ENTORNO (Suelo y Pared) ---
        \fill[pattern=north east lines] (-1.5, -0.2) rectangle (4, 0);
        \draw[thick] (-1.5, 0) -- (4, 0);
        \fill[pattern=north east lines] (-1.7, 0) rectangle (-1.5, 1.5);
        \draw[thick] (-1.5, 0) -- (-1.5, 1.5);

        % --- 2. RESORTE ---
        \draw[decoration={aspect=0.3, segment length=4pt, amplitude=3pt, coil}, decorate] 
            (-1.5, \hCar/2) -- (\xCar-\wCar/2, \hCar/2);
        \node[above] at ({(-1.5 + \xCar-\wCar/2)/2}, \hCar/2 + 0.1) {$k$};

        % --- 3. CARRO (Masa M) ---
        \draw[fill=gray!10, thick] (\xCar-\wCar/2, 0.1) rectangle (\xCar+\wCar/2, \hCar+0.1);
        \draw[fill=black] (\xCar-0.3, 0.1) circle (0.1);
        \draw[fill=black] (\xCar+0.3, 0.1) circle (0.1);
        \node[left] at (\xCar, \hCar/2+0.1) {$M$};
        \coordinate (O) at (\xCar, \hCar/2+0.1);
        \filldraw (O) circle (1.5pt) node[above right] {$O$};

        % --- 4. PÉNDULO (Masa m) ---
        \begin{scope}[shift={(O)}, rotate=\ang]
            \draw[ultra thick] (0,0) -- (0, -\L);
            \draw[line width=2.5pt,line cap= round] (0,0) -- (0, -\L);
            \draw[line width=1.5pt, gray,line cap= round] (0,0) -- (0, -\L);
            \shadedraw[ball color=gray!30] (0, -\L) circle (0.25) node[] {$m$};
            \node[right] at (0,-\L/2) {$L$};
        \end{scope}

        % --- 5. REFERENCIAS Y ÁNGULOS ---
        % Eje vertical de referencia
        \draw[dashed, gray] (O) -- (\xCar, -\L-0.5);
        % Ángulo theta
        \draw[->] (\xCar, {\hCar/2+0.1-0.8}) arc (270:{270+\ang}:0.8);
        \node at (\xCar+0.3, {\hCar/2+0.1-1.0}) {$\theta$};
        
        % Cota x(t)
        \draw[dashed, thin] (0, 1.2) -- (0, 0);
        \draw[|->, thick] (0, 1.0) -- (\xCar, 1.0) node[midway, above] {$x(t)$};

        % Vector Gravedad
        \draw[->, thick] (3.5, 1.5) -- (3.5, 0.8) node[below] {$\mathbf{g}$};

    \end{tikzpicture}
    \caption{Diagrama de un péndulo de masa $m$ acoplado a un soporte móvil de masa $M$ con restricción elástica.}
    \label{fig:pendulo_resorte_movil}
\end{figure}
% =========================================================================
% Capítulo 3: Cinética de sistemas de partículas
% Tema 01: Sistemas de partículas
% Problema 13: Péndulos acoplados con fuerza externa seguidora
% =========================================================================
\subsubsection{Péndulos acoplados con fuerza externa seguidora}

Considere el sistema de péndulos acoplados mostrado en la Figura \ref{fig:pendulos_acoplados}. El sistema está compuesto por una barra de longitud $l_1$ articulada en un punto fijo, a cuyo extremo inferior se articulan dos barras de longitudes $l_2$ y $l_3$. Todas las barras son rígidas e ideales (sin masa). En los extremos libres de las barras $l_2$ y $l_3$ se encuentran masas puntuales de valor $m$. Una fuerza externa $F$ actúa sobre la masa de la derecha, manteniendo siempre un ángulo fijo $\gamma$ con respecto a la línea de la barra $l_2$. Sabiendo que $\phi$ es el ángulo de la primera barra con la vertical y $\psi$ es el ángulo comprendido entre las dos barras inferiores, realice lo siguiente:

\begin{enumerate}
    \item Identifique las restricciones holónomas y no holónomas del sistema.
    \item Determine el número de grados de libertad.
    \item Obtenga las ecuaciones diferenciales del movimiento para las coordenadas generalizadas $\phi$ y $\psi$.
    \item Determine las fuerzas de restricción en las articulaciones.
    \item Analice si el sistema es conservativo y justifique su respuesta.
\end{enumerate}

\begin{figure}[!ht]
    \centering
    \begin{tikzpicture}[>=Latex, scale=1.4]
        % --- Parámetros ---
        \def\angPhi{-25}   % Ángulo phi (barra superior)
        \def\angPsi{80}   % Ángulo psi (apertura entre barras inferiores)
        \def\Lone{2.0}    % Largo l1
        \def\Ltwo{2}    % Largo l2
        \def\Lthree{1.5}  % Largo l3
        \def\angGamma{50} % Ángulo gamma de la fuerza

        % Coordenadas
        \coordinate (O) at (0,0); % Pivote superior
        \coordinate (J) at ({-90-\angPhi}:\Lone); % Articulación intermedia
        
        % Ángulos absolutos de las barras inferiores (asumiendo simetría respecto a la extensión de l1 para el dibujo)
        \def\angLeft{-90-\angPhi - \angPsi}
        \def\angRight{-90-\angPhi}
        
        \coordinate (M_L) at ($(J) + (\angLeft:\Lthree)$); % Masa izquierda
        \coordinate (M_R) at ($(J) + (\angRight:\Ltwo)$);  % Masa derecha

        % --- 1. Techo y Pivote ---
        \fill[pattern=north east lines] (-1, 0) rectangle (1, 0.2);
        \draw[thick] (-1, 0) -- (1, 0);
        \filldraw[white, draw=black] (O) circle (2pt);

        % --- 2. Barras ---
        % Barra l1
        \draw[line width=1.2pt] (O) -- (J) node[midway, right] {$l_1$};
        
        % Barras l2 y l3
        \draw[line width=1.2pt] (J) -- (M_L) node[midway, above left] {$l_3$};
        \draw[line width=1.2pt] (J) -- (M_R) node[midway, above right] {$l_2$};

        % Articulación J
        \filldraw[black] (J) circle (2pt);
                % --- 5. Fuerza F ---
        % Línea de extensión de l2
        \draw[dashed] (M_R) -- ($(M_R) + (\angRight:1.0)$);
        % Vector Fuerza F (a un ángulo gamma respecto a la extensión)
        \draw[->, thick] (M_R) -- ($(M_R) + ({\angRight-\angGamma}:1.2)$) node[right] {$F$};
        % Arco gamma
        \draw[->] ($(M_R) + (\angRight:0.6)$) arc (\angRight:{\angRight-\angGamma}:0.6);
        \node at ($(M_R) + ({\angRight-\angGamma/2}:0.8)$) {$\gamma$};

        % --- 3. Masas ---
        \shadedraw (M_L) circle (6pt) node[] {$m$};
        \shadedraw (M_R) circle (6pt) node[] {$m$};

        % --- 4. Ángulos ---
        % Referencia vertical para phi
        \draw[dashed] (O) -- (0, -1.5);
        \draw[->] (0, -0.8) arc (-90:{-90-\angPhi}:0.8) node[midway, below] {$\phi$};

        % Referencia extensión l1 y ángulo psi
        \draw[dashed, gray] (J) -- ($(J)!0.6!({-90-\angPhi}:4)$); % Extensión visual
        \draw[<-] ($(J) + (\angLeft:0.7)$) arc (\angLeft:\angRight:0.7) node[midway, below] {$\psi$};



        % --- 6. Gravedad ---
        \draw[->, thick] (-1.5, -1.0) -- (-1.5, -2.0) node[midway, left] {$g$};

    \end{tikzpicture}
    \caption{Sistema de péndulos acoplados con fuerza no conservativa.}
    \label{fig:pendulos_acoplados}
\end{figure}
% =========================================================================
% Capítulo 3: Cinética de sistemas de partículas
% Tema 01: Sistemas de partículas
% Problema 14: Mancuerna deslizante en pared y piso
% =========================================================================
\subsubsection{Mancuerna deslizante en pared y piso}

Una mancuerna está compuesta por dos partículas $A$ y $B$, cada una de masa $m$, unidas por una varilla rígida ideal (sin masa) de longitud $l$. El extremo superior $A$ de la mancuerna se desliza sin fricción a lo largo de una pared vertical fija, mientras que el extremo inferior $B$ se desliza sin fricción a lo largo de un piso horizontal. El ángulo $\theta$ se define entre la pared vertical y la varilla. El sistema se suelta desde el reposo en la posición vertical ($\theta=0$). Sabiendo que la gravedad actúa hacia abajo, realice lo siguiente:

\begin{enumerate}
    \item Defina los marcos de referencia y sistemas de coordenadas adecuados.
    \item Analice la cinemática del sistema para expresar la posición, velocidad y aceleración de las partículas en función del ángulo $\theta$.
    \item Analice la cinética del sistema para determinar la ecuación diferencial del movimiento (utilizando leyes de Newton y energía) y encuentre el ángulo $\theta$ en el cual la partícula $A$ pierde contacto con la pared vertical.
\end{enumerate}

\begin{figure}[!ht]
    \centering
    \begin{tikzpicture}[>=Latex, scale=1.5]
        % --- Parámetros ---
        \def\Lrod{3.5}    % Longitud de la barra l
        \def\ang{35}      % Ángulo theta con la vertical
        \def\wallH{3.5}   % Altura pared
        \def\floorW{3.0}  % Ancho piso

        % Coordenadas
        \coordinate (O) at (0,0);
        % Posiciones basadas en el ángulo theta con la vertical
        % A está en el eje Y: y = L*cos(theta)
        \coordinate (A) at (0, {\Lrod*cos(\ang)});
        % B está en el eje X: x = L*sin(theta)
        \coordinate (B) at ({\Lrod*sin(\ang)}, 0);

        % --- 1. Pared y Piso ---
        % Pared Vertical
        \draw[thick] (0, 0) -- (0, \wallH);
        \fill[pattern=north east lines] (-0.3, 0) rectangle (0, \wallH);
        \node[below left] at (O) {$O$};

        % Piso Horizontal
        \draw[thick] (0, 0) -- (\floorW, 0);
        \fill[pattern=north east lines] (0, -0.3) rectangle (\floorW, 0);

        % --- 2. Barra ---
        \draw[line width=3pt, gray!40] (A) -- (B);
        \draw[thick] (A) -- (B) node[midway, above right] {$l$};

        % --- 3. Masas ---
        % Masa A
        \shadedraw[ball color=gray!30] (A) circle (0.25);
        \node[left, xshift=-5pt] at (A) {$A$};
        \node[white, font=\footnotesize] at (A) {$m$};

        % Masa B
        \shadedraw[ball color=gray!30] (B) circle (0.25);
        \node[below, yshift=-5pt] at (B) {$B$};
        \node[white, font=\footnotesize] at (B) {$m$};

        % --- 4. Ángulo Theta ---
        \draw[dashed] (0, {\Lrod*cos(\ang)}) -- (0, {\Lrod*cos(\ang)-1.5}); % Referencia vertical desde A hacia abajo
        % O mejor, el ángulo suele medirse desde la pared en O o desde la vertical en A.
        % El dibujo original muestra el ángulo en la intersección de la barra con la vertical (pared).
        \draw[->, thick] (0, {\Lrod*cos(\ang)-0.8}) arc (270:{270+\ang}:0.8);
        \node at ({0.5*sin(\ang/2)}, {\Lrod*cos(\ang)-1.0}) {$\theta$};

        % --- 5. Gravedad ---
        \draw[->, thick] (2.0, 2.0) -- (2.0, 1.0) node[midway, right] {$\mathbf{g}$};

    \end{tikzpicture}
    \caption{Diagrama de cuerpo libre de una mancuerna deslizándose apoyada en una esquina.}
    \label{fig:mancuerna_esquina}
\end{figure}

\subsection*{Nota sobre soluciones}

Este capítulo aún no cuenta con soluciones desarrolladas para sus ejercicios.
Queda pendiente como trabajo futuro siguiendo el formato establecido en
\texttt{CONTRIBUTING.md}.
